% concatenated from neurips_2025.tex
\documentclass{article}


% if you need to pass options to natbib, use, e.g.:
\PassOptionsToPackage{numbers, compress,sort}{natbib}
% before loading neurips_2024


% ready for submission
\usepackage[final]{neurips_2025}
% \usepackage[preprint]{neurips_2024}

% to compile a preprint version, e.g., for submission to arXiv, add add the
% [preprint] option:
%     \usepackage[preprint]{neurips_2024}


% to compile a camera-ready version, add the [final] option, e.g.:
%     \usepackage[final]{neurips_2024}


% to avoid loading the natbib package, add option nonatbib:
%    \usepackage[nonatbib]{neurips_2024}


\usepackage[utf8]{inputenc} % allow utf-8 input
\usepackage[T1]{fontenc}    % use 8-bit T1 fonts
\usepackage{hyperref}       % hyperlinks
\usepackage{url}            % simple URL typesetting
\usepackage{booktabs}       % professional-quality tables
\usepackage{amsfonts}       % blackboard math symbols
\usepackage{nicefrac}       % compact symbols for 1/2, etc.
\usepackage{microtype}      % microtypography
\usepackage{xcolor}         % colors


\usepackage{graphicx} % Required for inserting images
\usepackage{amsmath}
\usepackage{amssymb}
\usepackage{amsthm}
\usepackage{enumitem}
\usepackage{geometry}
\usepackage{bm}
\usepackage{cleveref}
\usepackage{comment}

\newtheorem{theorem}{Theorem}
\newtheorem{proposition}{Proposition}
\newtheorem{corollary}{Corollary}
\newtheorem{lemma}{Lemma}
\newtheorem{remark}{Remark}
\newtheorem{definition}{Definition}
\newtheorem{example}{Example}

\newcommand{\cjp}[1]{{\footnotesize\textcolor{orange}{JC: #1}}}
\newcommand{\lt}[1]{{\footnotesize\textcolor{cyan!60!black}{\tt TL: #1}}}
\newcommand{\ql}[1]{{\footnotesize\textcolor{red}{QL: #1}}}

\title{A unified framework for establishing the universal approximation of transformer-type architectures}


% The \author macro works with any number of authors. There are two commands
% used to separate the names and addresses of multiple authors: \And and \AND.
%
% Using \And between authors leaves it to LaTeX to determine where to break the
% lines. Using \AND forces a line break at that point. So, if LaTeX puts 3 of 4
% authors names on the first line, and the last on the second line, try using
% \AND instead of \And before the third author name.


\author{%
  Jingpu Cheng\\
  % \thanks{Use footnote for providing further information
  %   about author (webpage, alternative address)---\emph{not} for acknowledging
  %   funding agencies.} 
      Department of Mathematics \\
   National University of Singapore\\
  \texttt{chengjingpu@u.nus.edu} \\
 \And
  Ting Lin \\
  School of Mathematical Sciences\\
  Peking University \\
  \texttt{lintingsms@pku.edu.cn} \\
  \And
  Zuowei Shen \\
    Department of Mathematics \\
   National University of Singapore\\
  \texttt{matzuows@nus.edu.sg} \\
    \And
  Qianxiao Li \\
  Department of Mathematics \\
 Institute for Functional Intelligent Materials \\
 National University of Singapore \\
  \texttt{qianxiao@nus.edu.sg} \\
}


\begin{document}


\maketitle

\begin{abstract}
    We investigate the universal approximation property (UAP) of transformer-type architectures, providing a unified theoretical framework that extends prior results on residual networks to models incorporating attention mechanisms. Our work identifies token distinguishability as a fundamental requirement for UAP and introduces a general sufficient condition that applies to a broad class of architectures. Leveraging an analyticity assumption on the attention layer, we can significantly simplify the verification of this condition, providing a non-constructive approach in establishing UAP for such architectures.
    We demonstrate the applicability of our framework by proving UAP for transformers with various attention mechanisms, including kernel-based and sparse ones. The corollaries of our results either generalize prior works or establish UAP for architectures not previously covered. Furthermore, our framework offers a principled foundation for designing novel transformer architectures with inherent UAP guarantees, including those with specific functional symmetries. We propose examples to illustrate these insights.
% Beyond recovering existing results, our theory enables the principled design of novel architectures; we propose examples with dynamic sparse attention that inherently satisfy UAP guarantees.
% By transforming UAP verification into a modular check on token-mixing dynamics, this work bridges theoretical expressivity and practical architecture innovation, offering a foundation for developing provably powerful transformers.
\end{abstract}

% \tableofcontents

\section{Introduction}

Transformers~\cite{vaswani2017attention} are a family of deep learning architectures that have 
achieved remarkable performance in natural language processing~\cite{radford2018improving,radford2019language,brown2020language}, computer vision~\cite{dosovitskiy2020image,carion2020end}, and other fields~\cite{jumper2021highly}. 
Given an input sequence of tokens, a transformer processes it through a deep composition of alternating attention and token-wise feedforward layers.  Besides the original softmax attention~\cite{vaswani2017attention}, a variety of different attention mechanisms have been proposed to enhance performance or computational efficiency, such as kernel-based attention~\cite{choromanski2020rethinking,tsai2019transformer,chen2021skyformer,katharopoulos2020transformers}, sparse attention~\cite{beltagy2020longformer,zaheer2020big,kitaev2020reformer}, and attention with low-rank structures~\cite{wang2020linformer,zhu2021long}.


% For an input $X_t\in\mathbb R^{d\times n}$ consisting of $n$ tokens with dimension $d$, a typical block in a transformer model builds the map $X_t\to X_{t+1}$ via the following two steps:
% \begin{align}
%     \label{eq:token_mixer}
%     X_{t+\frac{1}{2}} &= X_t + \operatorname{Atten}(X_t),\\
%     \label{eq:FFN}
%     X_{t+1} &= X_{t+\frac{1}{2}} + \operatorname{FFN}(X_{t+\frac{1}{2}}).
% \end{align}
% Here, $\operatorname{FFN}$ is applied token-wise on a token $X$ as:
% \begin{equation}
%     \operatorname{FFN}([X]_1,\dots,[X]_n) = (\operatorname{ffn}([X]_1),\dots,\operatorname{ffn}([X]_n)),
% \end{equation}
% where $\operatorname{ffn}:\mathbb{R}^d\to\mathbb{R}^d$ is typically a feedforward network. The operator $\operatorname{Atten}$ denotes the attention layer, which models inter-token interactions.
% In the original transformer~\cite{vaswani2017attention}, the attention layer is defined as
% \begin{equation}
%     \label{eq:atten}
%     \operatorname{Atten}(X_t)=\sum_{i=1}^N W_V^{t,i} X_t\,\operatorname{softmax}\Big((W_K^{t,i} X_t)^\top W_Q^{t,i} X_t\Big),
% \end{equation}
% with trainable parameters $W_V^{t,i},W_K^{t,i},W_Q^{t,i}\in\mathbb{R}^{d\times d}$ for the $i$th head in block $t$, and the softmax is applied column-wise. Recent years, a number of works have proposed alternative mechanisms to enhance performance or computational efficiency~\cite{beltagy2020longformer,zaheer2020big,choromanski2020rethinking,wang2020linformer,kitaev2020reformer,chen2021skyformer,guo2019star,tsai2019transformer}. For example, ... ...

% \ql{I think it is better to emphasize that transformer is not just one architecture,
% but a family of architectures based on different attention mechanisms.
% It may be also good to give some examples of variants of the attention mechanism,
% which allows you to segway to the next paragraph.}
% \ql{Also, need to emphasize the depth aspect.}

A natural theoretical question is: \emph{What is the expressive power of these architectures?} 
Previous studies have shown that transformers achieve the universal approximation property (UAP) via architecture-specific constructions, meaning they can approximate any continuous sequence-to-sequence function over compact domains~\cite{yun2020n,zaheer2020big,kajitsukatransformers} in certain measures. 
However, these results heavily rely on explicit, architecture-specific constructions, and a unified theoretical framework of deep transformer-type architectures remains elusive.
In particular, it is highly desirable to derive a 
verifiable condition that guarantees UAP of deep transformer-type architectures,
independent of specific architectural details such as the choice of the attention mechanism.
% However, these results depend heavily on explicit constructions. 
% Thus, a unified theoretical framework is highly desirable, particularly a general and verifiable condition guaranteeing UAP for transformers independent of specific architectural choices like attention mechanisms. 
Such a framework would allow greater flexibility in design without sacrificing expressivity.


% \ql{should limit discussion on depth aspect?}\cjp{have mentioned the depth part.}

Similar concerns have been addressed for fully connected deep residual networks (ResNets) using insights from control theory and dynamical systems~\cite{tabuada2022universal,ruiz2023neural,li2022deep,cheng2025interpolation}. By interpreting ResNets as control systems, recent studies~\cite{li2022deep,cheng2025interpolation} showed that deep ResNets with Lipschitz nonlinear activation functions possess UAP. However, extending this approach to transformers presents a challenge. Unlike ResNets, transformers apply identical feedforward transformations across tokens, without direct inter-token interactions. Hence, the attention mechanism must effectively capture token dependencies and propagate contextual information throughout the network.

To extend the UAP framework from ResNets to transformers,  we model each transformer block through two sequential operations\footnote{Here, we omit the layer normalization  for simplicity of the analysis.}:
\begin{equation}
    \begin{aligned}
    \label{eq:block}
    X_{t+\frac{1}{2}} &= X_t + \operatorname{Atten}(X_t),\\
    X_{t+1} &= X_{t+\frac{1}{2}} + \operatorname{FFN}(X_{t+\frac{1}{2}}),
    \end{aligned}
\end{equation}
where $X_t \in \mathbb{R}^{d\times n}$ represents $n$ tokens each of dimension $d$. The feedforward network ($\operatorname{FFN}$) acts independently on each token, while the attention layer ($\operatorname{Atten}$) explicitly models dependencies across tokens. 
A transformer model is then defined as the composition of such blocks. In many practical architectures, the attention layer is computed by some interactions between tokens, which has a computational complexity up to $\mathcal{O}(n^2d)$, while the feedforward layer has computation complexity of $\mathcal{O}(nd^2)$. This structure reduces computational complexity from $\mathcal{O}(n^2 d^2)$ (for using network with dense layers) to $\mathcal{O}(nd(n+d))$(can be even lower in many variants) by decomposing approximation tasks into simpler token-wise and token-mixing components. Such a decomposition not only enhances computational efficiency but also provides a novel perspective in the context of approximation theory. Therefore, it is of theoretical interest to understand \emph{how the combination of the token-wise and token-mixing operations contributes to the expressive power of models in handling sequential data.}

In this paper, we develop a general framework for the UAP analysis of transformers. Specifically, we provide abstract and verifiable conditions ensuring UAP, independent of specific architectural details. Our key contributions include:
% \ql{the following itemized parts seems too long - also don't change margins in any item}

\begin{itemize}
     \item We derive a general sufficient condition for
        transformer models to achieve UAP in the $L^p$ sense over compact sets (\Cref{thm:main}), requiring: (1) feedforward layers satisfying the conditions~\cite{cheng2025interpolation}, as stated in~\Cref{def:nonlinearity}, and (2) attention mechanisms producing distinct context-aware token representations across different inputs.
        % Specifically, we require the token-wise feedforward layers to satisfy the condition in~\cite{cheng2023interpolation}, as stated in~\Cref{thm:main}, which holds for most practical architectures. For the attention mechanism, we impose a key condition: it must be able to generate distinct, context-aware token representations across different data samples. 
        Notably, our framework  incorporates potential symmetry under token permutations in transformers, extending the analysis to non-transitive permutation groups beyond
        ~\cite{li2024deep}.

        \item For attention mechanisms that are analytic to their parameters, we reduce UAP verification to a two-sample condition (\Cref{prop:analytic}), simplifying practical validation compared to constructive approaches~\cite{yun2019transformers,yun2020n,zaheer2020big,kajitsukatransformers}.
        Moreover, we show conditions under which transformers with a fixed number of attention layers but arbitrarily many feedforward layers achieve UAP, generalizing the results on the memorization capability of transformers in~\cite{kajitsukatransformers,kim2022provable}.

        \item We apply our general framework to various transformer architectures to demonstrate its generality and applicability, including kernel-based~\cite{vaswani2017attention,choromanski2020rethinking,tsai2019transformer}, sparse~\cite{beltagy2020longformer,zaheer2020big,kitaev2020reformer}, and some other attention mechanisms~\cite{wang2020linformer,chen2021skyformer}. 
        % For kernel-based attention (formulated in~\Cref{sec:kernel_atten}), we derive a sufficient condition (\Cref{thm:kernel}) for UAP, which covers many existing architectures in~\cite{vaswani2017attention,choromanski2020rethinking,tsai2019transformer} and also indicates UAP for new architectures. 
        % We explicitly derive conditions that ensure UAP for these variants, extending and generalizing existing results.
        For kernel-based attention (formulated in~\Cref{sec:kernel_atten}), our result (\Cref{thm:kernel}) provides UAP guarantee for many existing architectures in previous works~\cite{vaswani2017attention,choromanski2020rethinking,tsai2019transformer} and also  for new architectures. 
        For sparse attention (e.g., architectures proposed in~\cite{beltagy2020longformer,zaheer2020big,kitaev2020reformer}), our result (\Cref{prop:sparse}) provides a UAP criterion which generalizes beyond the softmax attention and
        is free from technical assumptions on the sparse pattern.

        % we prove that sparse attention mechanisms satisfy UAP as long as their connection patterns maintain connectivity between tokens (\Cref{prop:sparse}). Compared to previous UAP results on sparse transformers~\cite{yun2020n,zaheer2020big}, our analysis is more general, extending beyond softmax attention, applicable to dynamical sparse patterns across layers, and does not assume additional technical details except for the connectivity condition. Moreover, our results indicate that a fixed number of sparse attention layers combined with arbitrarily many feedforward layers can achieve UAP, extending the results in~\cite{kajitsukatransformers} to sparse transformers.

        \item Our theoretical results also enable principled design of UAP-guaranteed architectures. We demonstrate this by proposing new transformer architectures with UAP guarantees, especially for attention mechanisms that preserve specific functional symmetry(\Cref{sec:new_atten}).

        % Our theoretical framework also provides insights on the design of new transformer architectures from the approximation side.
        % Specifically, we can propose new transformer architectures with inherent UAP guarantees based on our results. In particular, we explore new attention mechanisms capable of capturing specific functional symmetries with UAP guarantees, demonstrated using the dihedral group and cyclic group as examples~\Cref{sec:new_atten}.
\end{itemize}
We discuss in detail after our main results how they relate to the rest of the literature, and collect a more detailed review of related work in~\Cref{sec:related_work}.




% In the original transformer~\cite{vaswani2017attention}, the attention layer is defined as
% \begin{equation}
%     \label{eq:atten}
%     \operatorname{Atten}(X_t)=\sum_{i=1}^N W_V^{t,i} X_t\,\operatorname{softmax}\Big((W_K^{t,i} X_t)^\top W_Q^{t,i} X_t\Big),
% \end{equation}

% First, we derive a sufficient \ql{sufficient for? If you use this many times,
% perhaps introduce early what ``sufficient condition'' means thereafter.}
% condition on the kernel function for kernel-based attention, where attention scores are computed via kernel functions applied to query and key vectors~\cite{choromanski2020rethinking,tsai2019transformer}. This result covers many existing architectures, including the original transformer, Performer~\cite{choromanski2020rethinking}, and RBF-attention~\cite{tsai2019transformer}.
% Second, we investigate the UAP of sparse attention mechanisms (e.g., Longformer~\cite{beltagy2020longformer}, BigBird~\cite{zaheer2020big}, Reformer~\cite{kitaev2020reformer}), where each token attends only to a subset of all tokens.
% Our results not only recover existing UAP guarantees for sparse transformers but also address an open question from~\cite{child2019generating}: What connection patterns for attention suffice for universal approximation?
% We prove that sparse attention mechanisms satisfy UAP as long as their connection patterns maintain the strong graph connectivity.
% This aligns with the intuition that sparse attention must preserve an ``information highway'' between tokens across network depth.
% These findings highlight the generality and practicality of our framework, offering a unified theoretical lens for analyzing the UAP of transformer architectures.
% We believe our results provide a valuable theoretical foundation for novel transformer architectures.

% Notably, a similar idea termed ``contextual mapping'' has been employed in prior constructive proofs for original transformers and its sparse variants~\cite{zaheer2020big,yun2020n,kajitsukatransformers}.
% To our knowledge, this work is the first to unify this condition with the feedforward layer requirement into a general criterion for the UAP of transformers.

% \ql{paper should be self-contained, so briefly discuss what it is, or reference to
% the statement where this condition is given.},
% which holds for most practical architectures, regardless of the choice of activation functions.
% For the attention mechanism, we impose a key condition: it must be able to generate distinct, context-aware token representations across different data samples.
% Notably, a similar idea termed ``contextual mapping'' has been employed in prior constructive proofs for original transformers and its sparse variants~\cite{zaheer2020big,yun2020n,kajitsukatransformers}.
% To our knowledge, this work is the first to unify this condition with the feedforward layer requirement into a general criterion for the UAP of transformers.
% However, verifying this condition is still not easy in general, as it involves checking over arbitrarily many data samples.
% Another key contribution of our work is a more practical verification criterion based on an analytic assumption (satisfied by most existing architectures) on the attention mechanism, reducing the contextual mapping verification to only two different samples.
% Compared to prior explicit constructions, our results provide a general, verifiable UAP condition for transformers.

% To demonstrate the applicability of our framework, we apply
% our results for two classes of transformer generalizations.
% First, we derive a sufficient \ql{sufficient for? If you use this many times,
% perhaps introduce early what ``sufficient condition'' means thereafter.}
% condition on the kernel function for kernel-based attention, where attention scores are computed via kernel functions applied to query and key vectors~\cite{choromanski2020rethinking,tsai2019transformer}. This result covers many existing architectures, including the original transformer, Performer~\cite{choromanski2020rethinking}, and RBF-attention~\cite{tsai2019transformer}.
% Second, we investigate the UAP of sparse attention mechanisms (e.g., Longformer~\cite{beltagy2020longformer}, BigBird~\cite{zaheer2020big}, Reformer~\cite{kitaev2020reformer}), where each token attends only to a subset of all tokens.
% Our results not only recover existing UAP guarantees for sparse transformers but also address an open question from~\cite{child2019generating}: What connection patterns for attention suffice for universal approximation?
% We prove that sparse attention mechanisms satisfy UAP as long as their connection patterns maintain the strong graph connectivity.
% This aligns with the intuition that sparse attention must preserve an ``information highway'' between tokens across network depth.
% These findings highlight the generality and practicality of our framework, offering a unified theoretical lens for analyzing the UAP of transformer architectures.
% We believe our results provide a valuable theoretical foundation for novel transformer architectures.
% \ql{the last line is important to emphasize more, because ours can be used to
% design new architectures with guarantees. Do we have any examples?}

% \ql{Overall comment for this section: the length is fine for a ML paper.
% Maybe last two paragraphs can be made in itemized form.
% Some materials from sec 2 (e.g. the two-step architecture) can be moved here
% to consolidate space.}


\section{Problem formulation}
\label{sec:preliminaries}
In this section, we introduce the transformer-type architecture,
an abstraction of the standard transformer~\cite{vaswani2017attention} as a family of architectures composed of two repeating components: the token-mixing layers and token-wise map layers.
%  \lt{maybe we should say something about the abstraction?}
Then, we define the universal approximation property (UAP). Notably, we introduce the UAP under permutation equivariance for any subgroup $G$ of the symmetric group $S_n$ over tokens, which is a more general framework.
% \lt{in what sense? we should highlight that this is under the permutation equivariance. Therefore, to introduce a unified framework with such   symmetry is not easy}.

In the following, we use $X=([X]_1,\dots, [X]_n)\in\mathbb{R}^{d\times n}$ to denote one data sample consisting of $n$ tokens $[X]_1, \cdots, [X]_n$ of dimension $d$. We say that $X$ is in \emph{general position} if all of its tokens are distinct.
We will also use the notation $[n]:=\{1,\dots,n\}$ for any positive integer $n$.
% \lt{I think here $X$ stands for a series of indices. So maybe it is better to move the definition of the data space to here.}

\subsection{Transformer architecture}
We present a general formulation for the two-step architecture of transformers as described in~\eqref{eq:block}. The mapping $X_t \mapsto X_{t+1}$ can be abstracted as $(\operatorname{Id}+h)\circ (\operatorname{Id}+g)$. Here, $g$ generalizes the attention map $\operatorname{Atten}$ to a general token-mixing map, while $h$ generalizes the token-wise feedforward map $\operatorname{FFN}$, which applies to $X\in\mathbb{R}^{d\times n}$ as:
\begin{equation}
    \label{eq:tensor_type}
    h(X):=(\bar h ([X]_1),\dots,\bar h ([X]_n)),\quad \text{where } \bar h:\mathbb{R}^d\to\mathbb{R}^d.
\end{equation}

We denote by $\mathcal{G}$ a token-mixing family, consisting of functions mapping $\mathbb{R}^{d\times n}$ to $\mathbb{R}^{d\times n}$, to represent all possible choices of $g$ in a transformer. Recall that the attention layer in the original transformer~\cite{vaswani2017attention} is given by
\begin{equation}
    \label{eq:atten}
    \operatorname{Atten}(X_t)=\sum_{i=1}^N W_V^{t,i} X_t\,\operatorname{softmax}\Big((W_K^{t,i} X_t)^\top W_Q^{t,i} X_t\Big),
\end{equation}
with trainable parameters $W_V^{t,i},W_K^{t,i},W_Q^{t,i}\in\mathbb{R}^{d\times d}$ for the $i$-th head in block $t$, and the softmax is applied column-wise.
In this case, $\mathcal G$
 is precisely the family of functions defined by~\eqref{eq:atten} for all possible choices of $W_V^{t,i},W_K^{t,i},W_Q^{t,i}$.

Moreover, we consider
\begin{equation}
    \mathcal{H}^{\otimes n}:=\{X\mapsto (\bar h([X]_1),\dots,\bar h([X]_n)) \mid \bar h\in \mathcal{H}\},
\end{equation}
where $\mathcal{H}$ is a family of maps from $\mathbb{R}^d$ to $\mathbb{R}^d$, as the function family for the token-wise feedforward map $h$ in a transformer.
% \lt{I think we can adopt a tensor notation, say using $\mathcal H^{\otimes}$, or just using $\mathcal H$ (albeit a little abuse of notations)}
We define a transformer block, the generalization of~\eqref{eq:block}, to be a map in
\begin{equation}
    \mathcal{F}_{\mathcal{G},\mathcal{H}}:=\{(\operatorname{Id}+h)\circ (\operatorname{Id}+g) \mid g\in\mathcal{G},\, h\in\mathcal{H}^{\otimes n}\}.
\end{equation}
A transformer identified by $\mathcal G$ and $\mathcal H$ is then the composition of such blocks, i.e. a map in the set:
\begin{equation}
    \label{eq:transformer_hypothesis}
    \mathcal{T}_{\mathcal{G},\mathcal{H}}:=\{F_n\circ\cdots\circ F_1 \mid n\in\mathbb{N},\, F_i\in\mathcal{F}_{\mathcal{G}, \mathcal{H}}\}.
\end{equation}
% In practice, the attention layer (generalized as a token-mixing layer) in a transformer captures the dependencies between tokens, allowing language models to understand contextual relationships between words and vision models to integrate global spatial information across an image.
Notably, the feedforward layer can represent only tensor-type functions, i.e. functions of the form~\eqref{eq:tensor_type}.
The token-mixing mechanism extends this capability to more general functions by capturing the dependencies between tokens.


% A series of studies have demonstrated that deep neural networks can effectively approximate functions when operating on individual \(d\)-dimensional data points, typically with the number of parameters and computational complexity of \(\mathcal{O}(d^2)\) in each layer. \lt{cites?} However, when the input is a sequence \(X\in\mathbb{R}^{d\times n}\) of \(n\) tokens, directly applying a dense feedforward layer incurs a complexity of \(\mathcal{O}(n^2d^2)\). The two-step transformer-type architecture~~\eqref{eq:transformer_hypothesis} mitigates this issue by decoupling the approximation task into two components. Specifically, token-wise layers $\mathcal H$ can model single-token transformations on \(\mathbb{R}^d\), while the token-mixing layers are dedicated to modeling interactions among tokens. Notably, each of these components are responsible for one simpler task than modelling the full \(\mathbb{R}^{d\times n}\)-valued map.
% By the composition of these layers, this seperation reduces the layer-wise
% complexity to $\mathcal O(n^2+d^2)$ to approximate $\mathbb R^{d\times n}$-valued maps.
%  \lt{suggest a figure here.}
% The combination of inter-token and intra-token maps is not only practically advantageous, but also provides a new approximation scheme underexplored in classical approximation theory.
% Although previous works have provided explicit constructions to illustrate the approximation power of certain transformer architectures, our goal is to introduce a general framework that abstracts away from specific architectural details. This perspective paves the way for a deeper theoretical understanding of how the combination of the token-level and token-mixing operations contributes to the expressive power of models in handling sequential data.
% \ql{this paragraph reads like one in an introduction -- perhaps move some of these into the intro
% and make this part more concise.}

% \begin{remark}
%     For simplicity, normalization layers present in practical architectures are omitted.
% \end{remark}

\subsection{Universal approximation under permutation equivariance}
Let ${S}_n$ denote the symmetric group on $n$ elements and let $G\leq {S}_n$ be a subgroup. Then, $G$ has a natural group action over $\mathbb R^{d\times n}$ by permuting the $d$-dimensional tokens. A function $f:\mathbb{R}^{d\times n}\to\mathbb{R}^{d\times n}$ is said to be \emph{$G$-equivariant} if
\begin{equation}
    f(\sigma(X))=\sigma(f(X)),\quad \forall\,\sigma\in G,\; X\in\mathbb{R}^{d\times n}.
\end{equation}
The original transformer and many of its variants have some degree of permutation equivariance over tokens. For instance, kernel-based token mixers~\cite{vaswani2017attention,choromanski2020rethinking} typically have $G={S}_n$, whereas sliding-window attention~\cite{beltagy2020longformer} employs a binary group (identity and reflection), and some architectures~\cite{zaheer2020big,wang2020linformer} do not enforce any equivariance (i.e. $G=\{\operatorname{Id}\}$). If $\mathcal G$ consists of only $G$-equivariant functions, then $\mathcal T_{\mathcal G,\mathcal H}$ can approximate only $G$-equivariant target functions. This motivates the following definition:

\begin{definition}[$G$-UAP]\footnote{Notice that in~\Cref{def:G-UAP}, we do not require $\mathcal T_{\mathcal G,\mathcal H}$ to consist of only $G$-equivariant maps, but only that it can approximate $G$-equivariant functions. }
    \label{def:G-UAP}
    The transformer-type model with hypothesis space $\mathcal{T}_{\mathcal{G},\mathcal{H}}$ is said to have the \emph{$G$-universal approximation property (G-UAP)} in the $L^p$ sense ($1\le p<\infty$) if, for every continuous $G$-equivariant function $F:\mathbb R^{d\times n}\to\mathbb{R}^{d\times n}$, every compact set $K\subset\mathbb R^{d\times n}$, and every $\varepsilon>0$, there exists $\hat{F}\in\mathcal{T}_{\mathcal{G},\mathcal{H}}$ such that
    \begin{equation}
        \|\hat{F}-F\|_{L^p(K)}<\varepsilon.
    \end{equation}
\end{definition}


In applications, the equivariance restriction on the transformer is often addressed by introducing positional encoding~\cite{vaswani2017attention} on tokens. From a theoretical perspective, previous works~\cite{yun2019transformers,yun2020n,kajitsukatransformers,jiang2024approximation} have shown that if a family $\mathcal T_{\mathcal G,\mathcal H}$ has $G$-UAP for some $G$, then for any given compact set $K$, there exists an absolute positional encoding $\operatorname{Enc}:X\to X+E$, where $E$ is a fixed matrix, such that
\begin{equation}
    \mathcal T_{\mathcal G,\mathcal H}\circ \operatorname{Enc}:= \{ F\circ \operatorname{Enc}\mid F\in\mathcal F_{\mathcal G,\mathcal H} \}
\end{equation}
can approximate any continuous function on $K$ in the $L^p$ sense without symmetry constraints. Technically, this can be done by making the domains of each token position distinct. On the other hand, there are also applications where exact symmetry needs to be enforced, such as structure-to-property prediction in crystals
~\cite{xie2018crystal,chen2019graph,ren2020inverse,jiao2023crystal}.
Therefore, it is sufficient to consider the $G$-UAP, which can naturally extend to the general UAP while also covering cases where symmetry is considered. We will hereafter focus on the $G$-UAP.
% \lt{How about " in practice, there are two ways to resolve this issue. One is to consider the functions conforming to certain symmetry, the other is to use positional encodings. In this paper, we focus on the first approach.}

% \ql{I think one approach could be to say that with
% position encoding, any hyp space with G-equivariant UAP can be extended to
% have UAP in the general sense. This can be done by making the domains
% of each token position distinct. Thus, it is sufficient to consider the
% G-equivariant UAP, which also cover cases where exact symmetry needs
% to be enforced.}





In the literature, several works have studied the universal approximation of symmetric functions~\cite{cohen2016group,finzi2020generalizing,yarotsky2022universal,li2024deep}, often focusing on specific architectures and symmetric groups. Notably,~\cite{li2024deep} provides a general sufficient condition for the action of any transitive subgroup of $S_n$ on coordinates (1-dimensional tokens). In comparison, this work considers approximation under symmetry in a general setting, with group action over $d$-dimensional tokens instead of coordinates. Additionally, our results apply to non-transitive group cases, which are not covered in~\cite{li2024deep}. {In~\cite{agrachev2025generic}, the authors studied the ensemble controllability of control systems under symmetry, showing that systems that can interpolate arbitrarily many samples under symmetry are generic in a topology sense. However, this result does not tell us whether or not a given architecture has controllability. In comparison, our target is to provide a verifiable sufficient condition for UAP of specific architectures.}

Our analysis focuses on fixed-length sequence-to-sequence maps on compact subsets. This setting directly covers encoder-style tasks and many architectural variants,
which underlies many practical applications ranging from automatic speech recognition and visual sequence modeling to structure–property prediction in molecules and crystals~\cite{radford2019language,zheng2021rethinking,hatamizadeh2022unetr,jiao2023crystal,jumper2021highly}.
In parallel, there are also measure-theoretic formulations in the literature that treat inputs as probability measures, which can handle variable or even infinite context length under continuity/regularity assumptions~\cite{furuya2024transformers,geshkovski2024measure,geshkovski2025mathematical}, offering complementary insights to our results.

% \lt{Summarize the existing result with respect to symmetry. Can refer to our JMLR paper.}

% \subsubsection{Continuous-idealization}
% The continuous-time idealization of the transformer-type model is
% \begin{equation}
% \label{eq:ctn_tf}
%  \dot X(t)=\operatorname{FFN}_t(X(t))+\operatorname{Atten}_t(X(t)),\quad t\in [0,T].
% \end{equation}


% The control family of ~\eqref{eq:ctn_tf} is the following:
% \begin{equation}
% \label{eq:F}
%     \mathcal F:=\{\operatorname{FFN}+ \operatorname{Atten}\mid \operatorname{Atten}\in \mathcal G, \operatorname{FFN}\in \mathcal H\},
% \end{equation}
%     Here $\mathcal G$ and $\mathcal H$ are attention family and feedforward family, correspondingly.
% \begin{example}
%     A common setting for $\mathcal G$ and $\mathcal H $ is given by:
% \begin{equation}
% \label{eq:G}
%     \mathcal G=\{X\to W_V X\operatorname{softmax}((W_K X)^TW_Q X) \mid W_V, W_K, W_Q\in \mathbb R^{d\times d}\}
% \end{equation}
% \begin{equation}
% \label{eq:H}
%     \mathcal H = \{
%     X\to W_2(\sigma(W_1 X+b_1\mathbf{1}_n^T))+b_2\mathbf{1}_n^T \mid W_1\in\mathbb R^{m\times d}, W_2\in\mathbb R^{d\times m}, b_1\in\mathbb R^m, b_2\in\mathbb R^d
%     \}
% \end{equation}
% \end{example}


% In~\cite{cheng2023interpolation}, it is shown that under full affine invariance, the UIP and UAP hold for all non-linear Lipschitz control family.

% However, for the standard transformer model, there are two important restrictions in the control family $\mathcal F$ compared to the common affine invariant family:
% \begin{itemize}
%     \item \textbf{Permutation equivariance}: $\mathcal F$ consists of functions that are equivariant with respect to the token permutation, i.e. for any $f\in\mathcal F$ and $X\in\mathbb R^{d\times n}$, we have
%     \begin{equation}
%         f(XP)=f(X)P
%     \end{equation}
%     for any $n\times n$ permutation matrix $P$. This indicates that $\mathcal A_{\mathcal F}$(the flow map family) can at most approximate permutation equivariant function, and it cannot interpolate two samples whose tokens are permutation to each other but the targes are not.
    % \item\textbf{Restrictive affine invariance}: $\mathcal F$ possesses weaker affine invariance where the affine transformations are applied token-wisely. That is, the transformation matrices $W, A$ are restricted to be block diagonal(each block is $d\times d$), and $b$ is restricted to be the $n$ copy of a $d$-dimension vector.
% \end{itemize}

\section{Main results}
In this section, we establish a general sufficient condition for the UAP of transformer-type architectures. Since transformer architectures consist of token-wise feedforward layers and token-mixing attention layers, we first provide conditions for each component required for UAP.

For the feedforward family $\mathcal H^{\otimes n}$, we introduce the following definition:
\begin{definition}[Nonlinearity and affine-invariance for $\mathcal H$]
    \label{def:nonlinearity}
    We say a function family $\mathcal H$ (consisting of functions from $\mathbb{R}^d$ to $\mathbb{R}^d$) is nonlinear and affine-invariant, if 
    \begin{itemize}
        \item For any $h\in \mathcal{H}$ and any $W, A\in \mathbb{R}^{d\times d}, b\in \mathbb{R}^d$, the function $Wh(A\cdot -b)$ also belongs to ${\mathcal{H}}$;
        \item  $\mathcal{H}$ contains at least one non-affine Lipschitz function.
    \end{itemize}
\end{definition}
The nonlinearity and affine-invariance condition holds for almost all practical feedforward layers, independent of specific choices of activation functions and the width of the network. When $d\ge 2$, according to the main result in~\cite{cheng2025interpolation}, this condition ensures that the family
\begin{equation}
    h \in (\operatorname{Id}+\mathcal{H})^m = \{(\operatorname{Id}+h_m)\circ\cdots \circ (\operatorname{Id}+h_1) \mid h_1,\dots, h_m\in \mathcal{H} \}
 \end{equation}
can approximate any continuous function $f:\mathbb R^d\to\mathbb R^d$ in $L^p$ sense over compact set. Therefore, this condition guarantees that only the token-wise feedforward layer is able to generate complex features over a single token. 

However, an inherent limitation on the expressive power of feedforward layers is that they operate token-wise, meaning that they do not model any interactions between tokens. Considering this, we introduce the following definition for the attention family $\mathcal G$:

\begin{definition}[Token distinguishability for $\mathcal G$]
    \label{def:token_distinguishability}
    For a given group $G\le S_n$ and a set of samples $D:=\{X_i\}_{i=1}^N\subset \mathbb{R}^{d\times n}$ that are all in general position, we say a token-mixing family $\mathcal G$ can distinguish tokens in $D$ using $m$ layers under $G$-action, if there exists
    \begin{equation}
        g \in (\operatorname{Id}+\mathcal{G})^m = \{(\operatorname{Id}+g_m)\circ\cdots \circ (\operatorname{Id}+g_1) \mid g_1,\dots, g_m\in \mathcal{G} \}
     \end{equation}
     such that for any distinct $i,j\in[N]$ with $X_i$ and $X_j$ belonging to different orbits under the $G$-action (i.e., $X_i\neq \sigma(X_j)$ for all $\sigma\in G$), the tokens of $g(X_i)$ and $g(X_j)$ are all distinct.

     Moreover, we say $\mathcal G$ satisfies the \textbf{token distinguishability condition} under $G$-action, if for any finite set $D$, there exists $m$ such that $\mathcal G$ can distinguish tokens in $D$ using $m$ layers under $G$-action.

    % We say a token-mixing family $\mathcal G$, consisting of functions from $\mathbb R^{d\times n}$ to $\mathbb R^{d\times n}$, has token distinguishability with $G$-action if for any $N\in\mathbb{N}$ and any set of $N$ samples $\{X_i\}_{i=1}^N\subset \mathbb{R}^{d\times n}$ that are all in general position, there exist $m\in\mathbb{N}$ and
    % \begin{equation}
    %    g \in (\operatorname{Id}+\mathcal{G})^m = \{(\operatorname{Id}+g_m)\circ\cdots \circ (\operatorname{Id}+g_1) \mid g_1,\dots, g_m\in \mathcal{G} \}
    % \end{equation}
    % such that for any distinct $i,j\in[N]$ with $X_i$ and $X_j$ belonging to different orbits under the $G$-action (i.e., $X_i\neq \sigma(X_j)$ for all $\sigma\in G$), the tokens of $g(X_i)$ and $g(X_j)$ are all distinct.
    % That is, the token mixing layers can distinguish tokens across different samples.
\end{definition}

The token distinguishability condition ensures that token-mixing layers can model interactions between tokens by generating unique outputs for tokens in a finite set (up to $G$-action), enabling distinct in-context information for each token. This property is crucial for the expressive power of transformers, as illustrated below.

Consider a scenario where the token distinguishability condition fails: there exists a set $\Omega\in\mathbb R^{d\times n}$ with positive Lebesgue measure and some $i\in [n]$ such that $[g(X)]_i$ is constant over $\Omega$ for any $g\in (\operatorname{Id}+\mathcal G)^m$. Consequently, any $F\in\mathcal T_{\mathcal G,\mathcal H}$ is also constant over $\Omega$, leading to the failure of UAP. This example shows that if too many tokens are indistinguishable by the token-mixing mechanism (e.g., from a positive measure set), the transformer's expressive power becomes limited.

On the other hand, the token distinguishability condition is relatively mild, as it only demands the composition of token-mixing layers to distinguish tokens, rather than enforcing any precise relation. This condition is generally easy to satisfy, provided $\mathcal G$ includes sufficiently diverse maps that can effectively mix tokens.

In the following, we assume that $d\ge 2$ and the zero map is in $\mathcal G$. 
Based on~\Cref{def:nonlinearity,def:token_distinguishability},
% \ql{can you replace ``thanks to'' throughout?}
we can state our first main result on the UAP of transformers:
\begin{theorem}
    \label{thm:main}
    Suppose that $\mathcal H$ is nonlinear and affine-invariant~\Cref{def:nonlinearity}, and $\mathcal G$ satisfies the token distinguishability condition~\Cref{def:token_distinguishability}. Then, the family of transformers $\mathcal{T}_{\mathcal{G}, \mathcal{H}}$ satisfies the $G$-UAP~\Cref{def:G-UAP}.
    % \ql{hyperref the definitions for both parts}
\end{theorem}

\Cref{thm:main} provides a general condition for the UAP of transformers. However, directly verifying the token distinguishability condition is challenging since 
we need to check the condition arbitrarily many times.
% The following result significantly reduces the complexity in checking this condition, providing a more practical verification criterion under an analytic assumption on $\mathcal{G}$.
Therefore, we propose the following theorem, which greatly simplifies the procedure. 

\begin{theorem}
    \label{prop:analytic}
    We assume that $\mathcal{G}$ is parametrized by $
        \mathcal{G}=\{X\mapsto g(X; \theta) \mid \theta\in\Theta\subseteq \mathbb{R}^m\},$ where $\Theta$ is a connected open subset of $\mathbb{R}^m$, and for any fixed $X$, the mapping $\theta\mapsto g(X;\theta)$ is analytic. 
        Then, if $\mathcal{G}$ can distinguish tokens of any dataset $D$ with two elements (\Cref{def:token_distinguishability}) using finite many layers,
        % \ql{hyperref definition for distinguish tokens}
        then $\mathcal G$ satisfies the token distinguishability condition. 

    Moreover, if there exists a uniform $m$ such that with $m$ layers, $\mathcal G$ can distinguish tokens of dataset $D$ with $|D|=2$, then it can also do it for any finite dataset $D$ using $m$ layers.
    In this case, a deep model using only $m$ token-mixing layers and sufficiently many feedforward layers can achieve the UAP.
    % \ql{but still need enough ffn layers. Should mention this.}
\end{theorem}

The key insight in the proof of~\Cref{prop:analytic} is that if token distinguishability fails over a finite set, we can derive an equation in $\theta \in \Theta$ that is identically zero. By leveraging the property that the zero set of a non-trivial analytic function has measure zero, the equation can be reduced to the case of two elements, as detailed in~\Cref{sec:analytic_proof}. The use of the analytic property is straightforward but significantly simplifies the token distinguishability condition. 

Given expressive enough feedforward layers, \Cref{thm:main} highlights the role of token-mixing mechanisms in transformer architectures for UAP: generating distinct, context-aware token representations. This aligns with prior works~\cite{zaheer2020big,yun2020n,kajitsukatransformers}, which introduced ``contextual mapping'' to establish UAP for transformers. For instance,~\cite{kajitsukatransformers} defines ``contextual mapping'' as a function distinguishing tokens in a dataset $D$ (similar to $g$ in~\Cref{def:token_distinguishability}) without group actions. However, these works rely on explicit constructions, making verification complex and less generalizable. 
In contrast, \Cref{thm:main} is the first to our knowledge that formulates token distinguishability and feedforward layer conditions as a general, non-constructive criterion for UAP. 
There is no need to explicitly construct for UAP once the conditions are verified. 
Additionally, \Cref{prop:analytic} significantly simplifies the construction-based verification of token distinguishability, enabling broader applicability to diverse attention mechanisms, as shown in the examples in~\Cref{sec:app}. 
Furthermore, the uniformity of $m$ in \Cref{prop:analytic} also provides a convenient approach on the memorization capacity of attention layers studied in~\cite{kajitsukatransformers}.

% The uniformity of $m$ in \Cref{prop:analytic} also aids in analyzing the memorization capacity of attention layers, as studied in~\cite{kajitsukatransformers}.



% Given expressive enough feedforward layers, \Cref{thm:main} precisely
% identify the role of token-mixing mechanisms in transformer architectures with respect to UAP: the ability to generate distinct, context-aware representations for tokens. 
% This idea aligns with previous works~\cite{zaheer2020big,yun2020n,kajitsukatransformers}, where the concept of ``contextual mapping'' was introduced to establish the UAP of the original transformer and its sparse variants. For example, in~\cite{kajitsukatransformers}, a ``contextual mapping'' over a dataset $D$ is just defined as a function that can distinguish tokens in $D$ (just as $g$ in~\Cref{def:token_distinguishability}) with no group actions.
% However, these works all adopt constructive approaches to establish the UAP. To our knowledge, \Cref{thm:main} is the first to formulate the token distinguishability condition and conditions for feedforward layers as a general sufficient condition for the UAP of transformer-type architectures. Checking these conditions, there is no need to do explicit constructions for UAP. 
% Moreover, the existence of ``contextual mappings'' also relies on a constructive approach in previous works. The case-by-case construction makes the verification process be complex and difficult to generalize to different attention mechanisms.
% In comparison, \Cref{prop:analytic} significantly reduces the complexity in verifying the token distinguishability, allows the UAP results to be extended to a broader class of attention mechanisms easily, such as the examples we will present in~\Cref{sec:app}. 
% Furthermore, the uniformity of $m$ in \Cref{prop:analytic} also provides a convenient methods on studying the memorization capacity of attention layers studied in~\cite{kajitsukatransformers}.


% \subsection{General Sufficient Conditions for \texorpdfstring{$G$}{G}-UAP}
% \begin{theorem}
%     \label{thm:main}
%     Suppose $d\geq 2$, $0\in \mathcal{G}$ (the zero map). If the token-wise feature family $\mathcal{H}^{\otimes n}$ and the attention family $\mathcal{G}$ satisfy the following conditions:
%     \ql{Maybe some of these conditions above should be introduced as running assumptions,
%     so the main theorem only concerns the following two conditions.}
%     % \lt{Better to give each condition a name.}
%     \begin{enumerate}
%         \item \textbf{(Nonlinearity and affine-invariance for $\mathcal H$)} The function class $\mathcal{H}$ is affine-invariant in $\mathbb{R}^d$, meaning that for any $h\in \mathcal{H}$ and any $W, A\in \mathbb{R}^{d\times d}, b\in \mathbb{R}^d$, the function $Wh(A\cdot -b)$ also belongs to ${\mathcal{H}}$. Additionally, $\tilde{\mathcal{H}}$ contains at least one non-affine Lipschitz function.

%         \item \textbf{(Token distinguishability for $\mathcal G$)} For any $N\in\mathbb{N}$ and any set of $N$ samples $\{X_i\}_{i=1}^N\subset \mathbb{R}^{d\times n}$ that are all in general position, there exist $m\in\mathbb{N}$ and
%         \begin{equation}
%            g \in (\operatorname{Id}+\mathcal{G})^m = \{(\operatorname{Id}+g_m)\circ\cdots \circ (\operatorname{Id}+g_1) \mid g_1,\dots, g_m\in \mathcal{G} \}
%         \end{equation}
%         such that for any distinct $i,j\in[N]$ with $X_i$ and $X_j$ belonging to different orbits under the $G$-action (i.e., $X_i\neq \sigma(X_j)$ for all $\sigma\in G$), the tokens of $g(X_i)$ and $g(X_j)$ are all distinct.
%         That is, the token mixing layers can distinguish tokens across different samples.
%         % \ql{use ``general position'' since you have introduced it before}
%         % \lt{Maybe we can introduce $(Id + \mathcal G)^{\infty}$ as the union of all finite  $m$. }
%         % \lt{Good to include an illustration for this.}
%     \end{enumerate}
%     Then, $\mathcal{T}_{\mathcal{G}, \mathcal{H}}$ satisfies the $G$-UAP.
% \end{theorem}

% \ql{Sketch of proof should be given here.}

% \ql{A question: shouldn't condition 2 above require some minimality, i.e.
% it cannot distinguish *more* than $G$?}
% \cjp{If it can distinct more that $G$, then the transformer architecture can still appproximate any $G$-equivariant functions(it can actually approximate more). In our definition of $G$-UAP, we only require the transformer to approximate $G$-equivariant functions, but not restricting them to only represent $G$-equivariant functions.}

% Condition 1 in Theorem~\ref{thm:main} follows from~\cite{cheng2023interpolation} and ensures that compositions of feature maps can approximate any single-token function. This condition is satisfied by almost all practical feedforward layers, independent of specific choices of activation functions and the width of the network.

% Condition 2 ensures that the attention mechanism compensates for this limitation by enabling token mixing. Intuitively, it requires that the token mixer can produce unique outputs across all tokens in a finite sample set, allowing each token to incorporate distinct in-context information.

% \ql{I think some discussion of the consequence of this result should also be given here.
% Does this result say it is easy to have UAP for any mixing mechanism?
% why is distinguisability important? What if it fails, and do we have any
% simple examples to illustrate this?
% An alternative way to present it is to have two separate defintions,
% one for nonlinearity+affine invar and one for token distinguishability.
% After each of these definitions, discuss why having these properties
% may be important.
% Then, the theorem can be stated as a simple combination of these two requirements.
% and a proof sketch can be given before going on to the next part.}

% Although Theorem~\ref{thm:main} provides a general condition for the UAP of transformer-type architectures, verifying Condition 2 directly is challenging due to its dependence on an arbitrary number of samples. The following result provides a more practical verification criterion under an analytic assumption on $\mathcal{G}$.

% \begin{theorem}
%     \label{prop:analytic}
%     Suppose $\mathcal{G}$ is a parametric family, i.e.,
%     \begin{equation}
%         \mathcal{G}=\{X\mapsto g(X; \theta) \mid \theta\in\Theta\subseteq \mathbb{R}^m\}.
%     \end{equation}
%     Assume that $\Theta$ is a connected open subset of $\mathbb{R}^m$, and for any fixed $X$, the mapping $\theta\mapsto g(X;\theta)$ is analytic. Then, if $\mathcal{G}$ satisfies Condition 2 in Theorem~\ref{thm:main} for $N=2$ samples, it holds for any $N$ samples.

%     Moreover, if there exists a uniform $m$ such that $(\operatorname{Id}+\mathcal{G})^m$ can distinguish between any two samples \ql{should formally introduce this. Perhaps in a separate definition.
%     See commeont above}, then it can also do so for any $N$ samples. In this case, using only $m$ token-mixing layers and  suffices to achieve the UAP.
%     \ql{but still need enough ffn layers. Should mention this.}
%     % \lt{What can we obtain from the uniformity?}
% \end{theorem}

% Our results formally \ql{precisely?}
% identify the role of token-mixing mechanisms in transformer architectures with respect to UAP: the ability to generate distinct, context-aware representations for tokens. This idea aligns with previous works~\cite{zaheer2020big,yun2020n}, where the concept of ``contextual mapping'' was introduced to establish the UAP of the original transformer and its sparse variants.
% Compared to these prior results, where the contextual mapping condition was verified through explicit constructions, Proposition~\ref{prop:analytic} significantly simplifies the verification process by leveraging the analyticity of the attention mechanism.

% \ql{May need to explicitly explain what is ``contextual mapping'', and whether it coincides,
% or differs, from any concepts we introduced here.
% Make it clear to the reader what's new and what's known.}

% \begin{proposition}
%     Suppose $\mathcal G$ satisfies the conditions in Theorem~\ref{thm:main}. Then, the distinguishability condition 2 in Theorem~\ref{thm:main} holds for any $N$ and $N$ samples $\{X_i\}_{i=1}^N\subset \mathbb R^{n\times d}$ such that for each $X_i$, its $n$ tokens are all distinct. Specifically, there exists $m$ and $g\in (\operatorname{Id}+\mathcal G)^m$
% such that:
% For any $i,j\in[N]$ such that $X_i$ and $X_j$ are from different orbit of $G$, the tokens of $g(X_i)$ and $g(X_j)$ are all distinct.

% Moreover, when $\mathcal G$ is parametrized and analytic to all of its parameters, if there exists a uniform $m$ such that $(\operatorname{Id}+\mathcal G)^m$ can do the distinction for any 2 samples, then it can also do it for any $N$ tokens.
% \end{proposition}


% Consider any $N$ and $N$ samples $\{X_i\}_{i=1}^N\subset \mathbb R^{n\times d}$ such that for each $X_i$, its $n$ tokens are all distinct. There exists $m$ and
%         \begin{equation}
%             (\operatorname{Id}+\mathcal G)^m = \{(\operatorname{Id}+g_m)\circ\cdots \circ (\operatorname{Id}+g_1)\mid g_1,\cdots, g_m\in \mathcal G \}
%         \end{equation}
%         such that:

%  For any $i,j\in[N]$ such that $X_i$ and $X_j$ are from different orbit of $G$(i.e. there does not exists $g\in G$ such that $g(X_i)=X_j$), the tokens of $g(X_i)$ and $g(X_j)$ are all distinct.

\section{Applications to practical architectures}
\label{sec:app}

We demonstrate the generality and  applicability
of our UAP results by applying them to practical transformer architectures.
We first follow the kernel-based framework from~\cite{tsai2019transformer}, which provides a unified description for a series of attention mechanisms. Specifically, many attention variants proposed in prior work can be formulated as
\begin{equation}
    \label{eq:kernel_atten}
    [\operatorname{Atten}(X)]_i=\frac{\sum_{j\in\mathcal N(i)} k([W_QX]_i,[W_KX]_j)[W_VX]_j}{\sum_{j\in\mathcal N(i)} k([W_QX]_i,[W_KX]_j)}, \quad W_Q, W_K, W_V\in\mathbb{R}^{d\times d},
\end{equation}
where $k:\mathbb{R}^d\times \mathbb{R}^d\to \mathbb{R}^+$ is a positive kernel function, and $\mathcal N(i)\subset [n]$ denotes the set of indices that the $i$-th token attends to.
In the original transformer, the kernel function is defined as $k(x,y)=\exp(x^\top y)$, and $\mathcal{N}(i) = [n]$.

Under this framework, many transformer variants can be categorized into two types, to which we will apply our results::
\begin{itemize}
    \item \textbf{Kernel modification}: Replacing the kernel function $k$ to improve efficiency or performance. For example, using a kernel of the form $k(x,y)= \phi(x)^\top\phi(y)$ with a feature map $\phi:\mathbb{R}^d\to\mathbb{R}^m$ can significantly reduce computational cost when $m\ll n$.
    \item \textbf{Sparse attention}: For each $i$, restricting $\mathcal N(i)$ to a subset of $[n]$, reducing the number of tokens each token attends to. Here, we discuss in a general sense where $\mathcal N(i)$ can be dynamic across different layers, such as the sparse pattern in~\cite{zaheer2020big,kitaev2020reformer}. 
    % \lt{now we can alsy handle the dynamical case.}
\end{itemize}

% We now present the application of our results to these two categories.

\subsection{Kernel-based attention}
~\label{sec:kernel_atten}
We first consider the kernel modification case, where we assume $\mathcal N(i)=[n]$ for all $i$. The following result follows from \Cref{thm:main}:

\begin{corollary}
    \label{thm:kernel}
    % \ql{call this and similar ones corollaries?}
    Suppose the kernel function $k$ satisfies the following conditions:
    \begin{itemize}
        \item $k(\cdot,\cdot):\mathbb R^{d}\times \mathbb R^{d}\to \mathbb R^+$ is an analytic function.

        \item For any $x\in\mathbb{R}^d\setminus\{0\}$ and distinct points $y_1, y_2\in\mathbb{R}^d\setminus\{0\}$, for  almost all $W_K\in\mathbb{R}^{d\times d}$\footnote{means that the condition holds all the whole space except for a measure zero set.}, the following holds:
        \begin{equation}
 \lim_{t\to \infty}\frac{k(x, tW_Ky_1)}{k(x, tW_Ky_2)}= 0  \text{ or } +\infty.
        \end{equation}
        That is, for almost all given $W_K$, the kernel function $k$ can distinguish token representations by scaling the key vectors with a large factor.
        % For any $x\in\mathbb{R}^d\setminus\{0\}$ and distinct points $y_1, y_2\in\mathbb{R}^d\setminus\{0\}$, for almost all(in the Lebesgue measure sense) $W_K\in\mathbb{R}^{d\times d}$:
        % \begin{equation}
        %     \lim_{t\to \infty} \frac{1}{\sum_{j=1}^n k(x, tW_Ky_j)} \left(k(x, tW_Ky_1),\dots, k(x, tW_Ky_n)\right) = \delta_i, \quad \text{for some } i\in[n],
        % \end{equation}
        % where
        % $\delta_i= (0, 0, \dots, 1, \dots, 0, 0)\in\mathbb R^n$ is the $i$-th standard basis vector in $\mathbb{R}^n$. That is, scaling the key vectors $W_Ky_i$ by a large factor $t$ makes the attention scores concentrate on a certain token $i$.

    \end{itemize}
    Then, a transformer with kernel-based attention family $\mathcal{G}$ and feedforward family $\mathcal{H}$ satisfying the conditions in \Cref{thm:main} possesses the ${S}_n$-UAP. Moreover, using only one token-mixing layer and sufficiently many feedforward layers
    can achieve the UAP.
\end{corollary}

\Cref{thm:kernel} ensures the distinguishability condition in \Cref{thm:main} through the limiting behavior of the kernel function. This generalizes the idea from~\cite{yun2019transformers,yun2020n}, where softmax was used as an approximation of hardmax in explicit constructions. In comparison, our approach leverages analyticity, allowing the limit behavior to directly establish the distinguishability condition without further constructions.

Consequently, this result applies to various existing attention mechanisms.
In particular, the following kernels directly satisfy the condition in~\Cref{thm:kernel}:
\begin{itemize}
    \item $k(x,y)=\exp(x^\top y)$, used in the original transformer.
    \item $k(x,y)=\exp(-\gamma \|x-y\|_2^2)$ for $\gamma > 0$, the RBF kernel, explored in~\cite{tsai2019transformer}.
    \item $k(x,y)=\phi(x)^\top\phi(y)$, where
    \begin{equation}
        \phi(x)^{\top}=\exp \left(-\frac{1}{2}\|x\|^2\right)\left(\exp ({\omega}_1^{\top} x), \ldots, \exp ({\omega}_m^{\top} x)\right)\in\mathbb{R}^m,
    \end{equation}
    with $\omega_1,\dots, \omega_m\in\mathbb{R}^d$ being fixed weights drawn i.i.d. from a Gaussian distribution. This kernel is used in Performer~\cite{choromanski2020rethinking}, where Theorem~\ref{thm:kernel} holds almost surely.
\end{itemize}
Among these, the UAP for the original transformer and Performer have already been shown~\cite{yun2019transformers,alberti2023sumformer}. Our result recovers these results in our framework and relaxes the requirement on the architecture to achieve UAP: for original transformer, we do not need the bias in query vectors as in~\cite{yun2019transformers}; for Performer, we do not need additional hidden dimensions as in~\cite{alberti2023sumformer}.
To the best of our knowledge, the UAP for RBF kernel attention is new, demonstrating the generality of our approach. Moreover, we can easily propose other kernels satisfying the condition in Theorem~\ref{thm:kernel} but have not been studied in the literature, such as the following forms of $k(x,y):$
 \begin{itemize}
\item $k(x,y)=\exp(w^\top (x+y))$ for some $w\in\mathbb R^d\setminus\{0\}$;
        \item $k(x,y)=p(x-y)\tilde k(x,y)$, with $p$ being any positive polynomial function and $\tilde k$ being any kernel mentioned above.
 \end{itemize}


% \ql{Out of all of these, how many have UAP established and we are
% doing so more generally, and how many have unknown UAPs?
% Let's be clear about these things.}

% \ql{Shall we give some other kernels that satisfy the condition,
% but have not been studied in the literature?}

\Cref{thm:kernel}
also generalizes the results in~\cite{kajitsukatransformers} on the memory capacity of transformers, where they prove that for  transformers with dense softmax attention, one layer of attention is sufficient to achieve the UAP. Our result extends this to a broader class of kernel-based attention mechanisms.

\subsection{Sparse attention}
\label{sec:sparse_atten}

Prior works proposed sparse attention mechanisms to reduce the computational complexity of attention blocks~\cite{child2019generating,beltagy2020longformer,zaheer2020big,kitaev2020reformer, guo2019star,correia2019adaptively}. A common intuition for designing sparse patterns while retaining expressivity is ensuring connectivity, i.e., each token can attend to others via multiple “hops.” For instance, in sliding window attention~\cite{beltagy2020longformer}, where $\mathcal N(i)=\{j\in[n]\mid |j-i|\le w\}$ with $w\ll n$, long-range interactions are achieved indirectly via multiple attention layers. In the following, we formalize this intuition and provide a general UAP condition for sparse attention transformers as a direct consequence of \Cref{thm:main}.

Denote $P([n])$ as the power set of $[n]$, i.e. the set of subsets of $[n]$.
For a given function $\mathcal N:[n]\to P([n])$, we define
$\mathcal G_{\mathcal N}$ as the family of maps  from $\mathbb R^{d\times n}\to\mathbb R^{d\times n}$ defined by~\eqref{eq:kernel_atten} associated with the sparsity pattern $\mathcal N$. 
We define the adjacency matrix of $\mathcal N$ as an $n\times n$ matrix $A_{\mathcal N}$ with $A_{\mathcal N}(i,j)=1$ if $j\in \mathcal N(i)$ and $A_{\mathcal N}(i,j)=0$ otherwise.


We also define
\begin{equation}
    \operatorname{Aut}(\mathcal N):=\{\sigma\in S_N\mid j\in \mathcal N(i)\Leftrightarrow \sigma(j)\in \mathcal N(\sigma(i))\}
\end{equation}
as the permutations that keep the structure of $\mathcal N$ invariant.

Let $\Phi:=(\mathcal N_1,\mathcal N_2,\cdots, )$ be a sequence of sparsity patterns. We define the sparse transformer family associated with $\Phi$ as:
\begin{equation}
    \mathcal T_{\mathcal H}^\Phi:=\{(\mathrm{Id}+h_n)\circ(\mathrm{Id}+g_n)\circ \cdots \circ(\mathrm{Id}+h_1)\circ(\mathrm{Id}+g_1)\mid n\in \mathbb N_+, h_i\in \mathcal H^{\otimes n}, g_i\in\mathcal G_{\mathcal N_i}, \text{ for } i\in [n]\}.
\end{equation}
Such a definition formulates transformers with dynamic sparse attention patterns. 
% \ql{I feel that the following should be a definition}

\begin{definition}
% Suppose that there exists $m$ such that for any $i\neq j\in [n]$, there exists $1\le r_1<r_2<\cdots <r_k\le m$ and $l_1, l_2,\cdots, l_k\in [n]$ with $l_1=i, l_k=j$ such that
% \begin{equation}
%     l_{p+1}\in \mathcal N_{r_p}(l_p), \quad \text{for } p\in [k-1],
% \end{equation}

We call the sparsity pattern $\Phi$ to be connected within $m$ layers, if for any $i\neq j\in [n]$, there exists a sequence $1\le r_1< r_2<\cdots <r_k\le m$ such that
\begin{equation}
    A_{\mathcal N_{r_k}}A_{\mathcal N_{r_{k-1}}}\cdots A_{\mathcal N_{r_1}}(i,j)>0.
\end{equation}
That is, any token can reach any other token through a subsequence of the $m$ sparse attention layers.
\end{definition}
% \ql{is this more easily stated as a connectivity condition on some graph, i.e.
% that corresponding to the product of adjancency matrices of the $\mathcal N_i$?}

 Also, let $\mathcal H$ be a family of token-wise feedforward layers satisfying the condition in~\Cref{thm:main}, and $k$ be a kernel function satisfying the condition in~\Cref{thm:kernel}.
Then, we have the following result:
\begin{corollary}
    % \cjp{Or we can add condition that there are some $\mathcal N_i$ with $\operatorname{Aut}(\mathcal N_i)=\{Id\}$, and the result is UAP without symmetric restriction.}
    % \ql{I think it is not necessary to discuss non-symmetric cases,
    % since they can be given as corollary by just adding position encoding.
    % if you must, then better to do it after this result.}
    \label{prop:sparse}

    Suppose that $\Phi$ is connected within $m$ layers.
    Then, $\mathcal T_{\mathcal F}^\Phi$ possesses the $G$-UAP, where
    \begin{equation}
        \label{eq:sparse_group}
        G=\bigcap_{i=1}^\infty \operatorname{Aut}(\mathcal N_i).
    \end{equation}
    Moreover, transformer with only $m$ layers of attention associated with the sparsity patterns $\mathcal N_1,\cdots, \mathcal N_m$ in $\Phi$ and sufficient number of token-wise feedforward layers can achieve the UAP.
\end{corollary}


% \begin{proposition}
%     \label{prop:sparse}
%     Let $\mathcal H$ be a family of token-wise feedforward layers satisfying the condition in Theorem~\ref{thm:main}.
%     Let $k$ be a kernel function satisfying the condition in Theorem~\ref{thm:kernel}, and let $\mathcal{G}$ be the attention family represented by~~\eqref{eq:kernel_atten}, with $\mathcal N(i)\subseteq [n]$ for all $i$. Define the directed graph $\Gamma=(V, E)$ associated with $\mathcal{G}$ as
%     \begin{equation}
%         V=\{1,\dots, n\}, \quad i\to j\in E \text{ if and only if } j\in\mathcal N(i).
%     \end{equation}
%     Then, if $\Gamma$ is \emph{strongly connected},
%     we have $\mathcal T_{\mathcal G,\mathcal H}$
%     possesses the $\operatorname{Aut}(\Gamma)$-UAP, where $\operatorname{Aut}(\Gamma)$ denotes the automorphism group of $\Gamma$. Moreover, the minimal number of token-mixing layers used in the transformer to achieve the UAP is no more than the diameter of $\Gamma$.
% \end{proposition}

\Cref{prop:sparse} provides a rigorous justification that the heuristic in keeping connectivity in the attention layers is also sufficient for UAP.
% \ql{I think the heuristic should be that it is necessary, but it being sufficient
% should be new?}
Results from graph theory indicate that when $n$ is large, a random sparse pattern $\mathcal N$ has a trivial automorphism group with probability approaching 1~\cite{erdos1963asymmetric}. This fact indicates that with the guarantee of connectivity, most of the sparse attention patterns
allow the UAP without symmetric restriction even in the absence of positional encodings.

% \cjp{Our result can also be extended to transformers with dynamic sparse attention patterns across layers. However, may need to introduce additional notations and change the definition of $\mathcal G$ and $\mathcal T_{\mathcal G, H}$ for a rigorous statement. Maybe we should move it to technical section.}  \lt{I think we can only gives a definition and a statement but without any lengthy technicalities. After all, I think the strategy of proofs work here.}

\Cref{prop:sparse} can cover many existing sparse attention mechanisms, including the following:
\begin{itemize}
    \item the periodic pattern switching between ``fixed attention'' and ``strided attention'' in~\cite{child2019generating};
    \item the sliding window attention with/without global seeds (tokens connect to all others) in~\cite{beltagy2020longformer};
    \item the star-shape attention in~\cite{guo2019star}, where one token attends to all others and the others connect in a circle;
    \item BigBird, a mixture of sliding windows, global seeds and random connections in~\cite{zaheer2020big}.
\end{itemize}

% Prior results on UAP for sparse softmax transformers~\cite{yun2020n,zaheer2020big} indicate that when the sparsified attention retains specific connection structures, the transformer still satisfies UAP. However, there are cases where their results do not apply, and the general condition for UAP of sparse transformers remains unknown.
% The following result, as a corollary of our main theorem, provides a general condition for UAP of transformers with fixed sparse attention mechanisms.

The UAP of transformers with sparse softmax attention have also been studied in~\cite{yun2020n,zaheer2020big} via a constructive approach.
Compared to their results, \Cref{prop:sparse} has several advantages. First, our results for UAP do not require other technical conditions, such as the periodicity of the sparse patterns and the existence of Hamiltonian path in~\cite{yun2020n}, or each sparse pattern contains a star sub-structure in~\cite{zaheer2020big}, other than the connectivity of the graph. For example, if the connection mode $\Phi$ is not periodic (e.g. there are different random patterns across layers), and some of the $\mathcal N_i$ do not contain a star graph mode, our result can still be applied as long as the connectivity is kept, while the results in~\cite{yun2020n,zaheer2020big} may not be applicable.

In addition, our results can be applied to all kernels that satisfy the condition in \Cref{thm:kernel}, which generalizes the results based on explicit construction using softmax attention. Moreover, we identify the number of token-mixing layers required to achieve UAP as the minimal number of ``hops'' for each token to attend to all other tokens.
In contrast, the results in~\cite{yun2020n,zaheer2020big} use unbounded number of layers to achieve UAP.
In this regard, our result can also be viewed as a generalization of the results on the minimal number of attention layers for UAP in~\cite{kajitsukatransformers} to sparse transformers.

% \ql{are there any sparsity patterns
% not considered previously that we can show have UAP?}

\subsection{Other attention mechanisms}
% \cjp{add a new subsection for other attention mechanisms.}
Our framework can also be conveniently applied to many other variants of attention mechanisms that cannot be covered by~\eqref{eq:kernel_atten}. For example, the token distinguishability condition can be verified for the following architectures using similar method as in~\Cref{thm:kernel}:
\begin{itemize}
    \item Linformer~\cite{wang2020linformer}, where the attention layer is defined as
    \begin{equation}
        \label{eq:linformer}
        \operatorname{Atten}(X)=W_V^{t} X {F}\operatorname{softmax}((W_K^{t} X{E})^\top W_Q^{t} X),
\end{equation}
where $E, F\in \mathbb R^{n\times k}$ with $1\le k\ll n$ are two trainable projection matrices. This variant of attention reduces the complexity of attention from $\mathcal O(n^2)$ to $\mathcal O(nk)$.
    \item Kernelized attention used in SkyFormer~\cite{chen2021skyformer}, where the attention mechanism is given by:
    \begin{equation}
        \label{eq:skyformer}
        [\operatorname{Atten}(X)]_i= \sum_{j=1}^n \exp{\left(-\frac{1}{2}\|[W_QX]_i-[W_KX]_j\|^2\right)}W_VX_j.
    \end{equation} 
\end{itemize}
\begin{corollary}
    We have that: (i)
LinFormer satisfies the UAP without symmetric restriction; (ii) SkyFormer satisfies the $S_n$-UAP.
\end{corollary}
The proofs are provided in~\Cref{sec:other_transformer} using~\Cref{prop:analytic}, similar to the proof of ~\Cref{thm:kernel}.

\subsection{New attention mechanisms from the approximation analysis}
\label{sec:new_atten}

% \ql{This part is important, but we should discuss it in a way that is interesting
% without the need for extensive nuemrics.}
% \ql{We should discuss some interesting aspects of this construction.
% For example, can we explicitly construct one that is invariant to some simple subgroup of $S_n$?
% Perhaps use crystal gen as an example? We should discuss this point.}

Our results also provide insights into designing new transformer architectures with inherent UAP guarantees. In particular, our framework inspires the design of architectures with UAP under specific symmetries. In this section, we present examples of such designs to illustrate these insights.

\subsubsection{New attention mechanism with bias term}

We propose a new architecture that naturally satisfies the conditions in \Cref{thm:main} and \Cref{prop:sparse}. We consider the following attention mechanism with bias term:
\begin{equation}
    \label{eq:new_atten}
    [\operatorname{Atten}(X)]_i=[X]_i+\sum_{j\in\mathcal N(i)} a \alpha(W[X]_j-b),
\end{equation}
where $a\in \mathbb R, W\in\mathbb R^d$ and $b\in\mathbb R$ are learnable parameters.
Assume that $\alpha$ is of polynomial growth, i.e. there exists $M$ and $N$ such that $|\alpha(x)|\le M(1+|x|^N)$ for all $x\in\mathbb R$.

Then, the result in~\Cref{prop:sparse} still holds, if we replace the attention mechanism in~\eqref{eq:kernel_atten} with~\eqref{eq:new_atten}. See~\Cref{sec:proof_new_atten} for the formal statement and proof.
    % Moreover, transformer with only $m$ layers of attention associated with the sparsity patterns $\mathcal N_1,\mathcal N_2,\cdots, \mathcal N_m$ in $\Phi$ and sufficient number of token-wise feedforward layers can achieve the UAP. \lt{not understand.} \cjp{Just want to say that the same results for sparse attention hold if we replace the attention mechanism with the new one.}

% We have the following proposition:
% \begin{proposition}
%     \label{prop:new_atten}
%     If $\alpha$ is analytic and not a polynomial. If $\Gamma$ is strongly connected, and the feedforward family $\mathcal{H}^{\otimes n}$ satisfying the conditions in Theorem~\ref{thm:main}.

%     Then, $\mathcal T_{\mathcal G,\mathcal H}$ possesses the $\operatorname{Aut}(\Gamma)$-UAP, where $\operatorname{Aut}(\Gamma)$ denotes the automorphism group of $\Gamma$.
%     \ql{Define aut somewhere}
%     Moreover, the minimal number of token-mixing layers required to achieve the UAP is no more than the diameter of $\Gamma$ \ql{Define diam somewhere}.
% \end{proposition}

% That is, the same results as Proposition~\ref{prop:sparse} can be achieved by using this simpler token-mixing map for any non-polynomial function.
\subsubsection{Transformer with UAP under specific symmetry}

% In many application scenarios, one needs architectures with specific symmetric restrictions.
% From this aspect, our general framework also provides a new perspective on the design of architectures with UAP guarantee under symmetric restrictions. 
% Specifically, for a given permutation group $G\le S_n$, we can design $G$-equivariant token-mixing layers that satisfy the token distinguishability under $G$-action. Then, according to~\Cref{thm:main}, a transformer model with such token-mixing layers and a feedforward family $\mathcal H$ with nonlinearity and affine invariance have $G$-UAP.
% This approach significantly reduces the considerations on expressive power in designing the architectures, since
% only the token distinguishability condition is required for the token-mixing layer. 

In many applications, architectures with specific symmetric restrictions are required. Our framework also offers a new perspective on designing such architectures with UAP guarantees. For a given permutation group $G\le S_n$, we can design $G$-equivariant token-mixing layers that satisfy token distinguishability under $G$-action. By~\Cref{thm:main}, a transformer with such token-mixing layers and a feedforward family $\mathcal H$ with nonlinearity and affine invariance achieves $G$-UAP. This simplifies the design process, as only token distinguishability is required for the token-mixing layer.

For some subgroups $G$ of $S_n$, the design for token-mixing layers can be very simple. Here, we use the example of $G=D_n$, the dihedral group of order $2n$, and the cyclic group $C_n$ of order $n$ to demonstrate. We identify $D_n\le S_n$ as the group generated by the cycle $\rho:=(1,2,\cdots,n)$ and the reflection $\sigma:=(1,n)(2,n-1)\cdots$. 
$C_n$ is the cyclic group generated by $\rho$.
$D_n$ corresponds to the symmetry of a regular $n$-gon, relevant in applications like molecular structure~\cite{bunker2006molecular,kittel2018introduction,gilmer2017neural}. Symmetry under $C_n$ applies to modeling periodic data, such as periodic time series~\cite{he2019stcnn,fan2022depts} and classifying periodic variable stars in cosmology~\cite{zhang2021classification}.


% \ql{May need to be more specific than cosmology. Also, isn't this also related to convolution network applications, say 
% temporal convolution nets (Haotian's work)?}

For $D_n$, we provide the designs of token-mixing layers with token distinguishability.
\paragraph{Architecture with $D_n$-symmetry}
 Choose the token-mixing layers defined in~\eqref{eq:kernel_atten} or in~\eqref{eq:new_atten} with the sparsity pattern 
    \begin{equation}
        \mathcal N(i):=\{(i+j)\operatorname{mod} n \mid j=-w,\cdots, 0,\cdots, w\},
    \end{equation}
    where $w\le\lfloor \frac{n-1}{2}\rfloor-1$ is an integer. Here, we assume that the kernel $k$ in~\eqref{eq:kernel_atten} and the function $\alpha$ in~\eqref{eq:new_atten} satisfy the conditions in~\Cref{thm:kernel,thm:new_atten}, respectively.

For $C_n$, similar designs as above also work. However, we can use a simpler one based on convolution:
\paragraph{Architecture with $C_n$-symmetry:}  Define the token-mixing layer via column-wise convolution:
    \begin{equation}
        \label{eq:conv}
        \operatorname{Atten}(X)= \psi * X, \quad \text{where }
        [\psi * X]_i = \sum_{j=0}^{l} \psi_{j}\, [X]_{(\ell + j) \bmod n}, \quad i = 0,1,\dots,n-1.
    \end{equation}
    with a trainable kernel $\psi=[\psi_{0},\cdots, \psi_l]\in\mathbb R^{l+1}$ for some integer $l\ge 1$.

% \ql{describe the difference between this and (temporal) convolutional networks?}
This design can be viewed as an adaptation of the temporal convolutional network~\cite{lea2016temporal,jiang2023forward}, treating the input sequence as a circular structure. The following statement holds, with proof in~\Cref{sec:proof_new_atten}:
\begin{corollary}
    The architecture with $D_n$-symmetry and $C_n$-symmetry defined above satisfies the token distinguishability condition under the action of $D_n$ and $C_n$, respectively. Moreover, the transformer with such token-mixing mechanisms and a non-linear affine-invariant family $\mathcal H$ possesses the $D_n$-UAP or $C_n$-UAP, respectively.
\end{corollary}
% We can easily prove that both architectures satisfy the token-distinguishability condition under their corresponding groups. Moreover, it is clear that the functions represented by them are indeed equivariant. Therefore, the transformer with such token-mixing maps and a feedforward family $\mathcal H$ satisfying the conditions in~\Cref{thm:main} or~\Cref{thm:new_atten} are architectures with desied symmetry and UAP guarantee. See for more details.

Architectures incorporating symmetry through convolutional layers have been studied in the literature ~\cite{{cohen2016group,cohen2018spherical,weiler20183d,zhang2021classification,lin2022universal}}. In particular,~\cite{zhang2021classification} also proposed a convolutional structure for representing $C_n$-invariant functions, although it does not inherently guarantee UAP.
In contrast, our proposed architecture naturally satisfies the $C_n$-UAP property according to \Cref{thm:main}.
Moreover, by choosing specific sparse mode $\mathcal N$, we can generalize the architecture for $D_n$ symmetry to other permutation subgroups of $S_n$ that can be identified as the automorphism group of an order-$n$ directed graph. 
For more general permutation groups, our framework can still be applied if one can find token-mixing layer satisfying the token distinguishability condition under the action of the group. The method proposed in~\cite{romero2020group,hutchinson2021lietransformer} may be helpful in identifying such layers. We believe our framework provides a new perspective on the design of equivariant/invariant architectures with UAP guarantees.

% \ql{not precise enough, which ones? Don't criticize, but objectively say which do not have approximation guarantees}



% As a byproduct, Proposition~\ref{prop:sparse} may also have implications on the architecture for permutation equivariant maps for a given group. In general, identifying a permutation group as the automorphism group requires about the same number of nodes as the group size. However, there still exists some groups that can be identified by a small graph. In this case, the transformer architecture may be a candidate for efficient representation of the permutation equivariant maps.



% The key idea is that they propose to use sparse connections in the computation of pair-wise attention, i.e. each token compute the softmax attention only with a subset of all the tokens. Spacifically, the sparsified attention machanism can be written as
%     \begin{equation}
%         \label{eq:sparsified_atten}
%         \operatorname{Atten}(X)_i=X_i+ \sum_{j\in \mathcal N(i)} \left(\frac{e^{\left\langle W_Q X_{s^\prime}, W_K X_j\right\rangle}}{\sum_{l\in \mathcal N(i)}  e^{\left\langle W_Q X_{s^\prime}, W_K X_{l}\right\rangle}} W_V X_j\right)
%     \end{equation}
%     where $\mathcal N(i)\subseteq [n]$ is the subset where the $i$-th token doing attention with.
\section{Conclusion}
In this paper, we investigate the universal approximation property (UAP) of general transformer architectures within a unified framework. Our main results, \Cref{thm:main} and \Cref{prop:analytic}, provide general and verifiable conditions for establishing UAP across a range of attention-based architectures, avoiding complex constructions as in previous works. This generality is demonstrated in \Cref{sec:app}, where we apply our framework to various attention types. Moreover, our results offer guidance for designing new attention mechanisms with UAP guarantees, as illustrated in \Cref{sec:new_atten}.
We also acknowledge certain limitations of this work. First, normalization layers commonly used in practice are not considered, and extending our analysis to incorporate them would be valuable. Second, some architectures, such as those in~\cite{katharopoulos2020transformers}, do not satisfy the analyticity assumption in \Cref{prop:analytic}. Although the condition in \Cref{thm:main} remains verifiable for such architectures, it remains unclear whether our results on the required number of token-mixing layers for UAP still hold.
Moreover, as our results offer non-constructive yet verifiable criteria for UAP—abstracting away the specific forms of token-mixing and token-wise modules—they do not yield quantitative insight into the relative contributions of each architectural component.
A systematic, quantitative characterization of how individual mechanisms (e.g., multi-head attention, mixture-of-experts, low-rank projections) affect approximation efficiency remains an important direction for future work.




\begin{comment}
\section{Proof Ideas}


\cjp{Seems no enough space for the proof ideas. Is it still necessary to discuss a little bit in the main text?}

\ql{You can add a few sentences which describe this section. 2-3 sentences about the things you want readers
to pay attention to, that represents the key theoretical insights or new techniques, then move this to the start of the proof appendices.}


\subsection{Proof Outline for \Cref{thm:main} and \Cref{prop:sparse}}

\ql{I think this part should highlight the technical novelties of the proof,
so it is better to discuss FIRST the new ideas, e.g. starting from the last paragraph.}

Our proof framework is based on~\cite{li2022deep,cheng2025interpolation}.
\ql{update these references to the published versions}
Following their approach, given a compact set $K\subset \mathbb{R}^{d\times n}$,
the function class $\mathcal{T}_{\mathcal{G}, \mathcal{H}}$ satisfies the UAP in the $L^p(K)$ norm ($1\leq p<\infty$) if the following two conditions hold:
\begin{itemize}
    \item \textbf{Interpolation Property}: For any $\varepsilon>0$ and any finite dataset $\{(X_i, Y_i)\}_{i=1}^N\subset K\times \mathbb{R}^{d\times n}$, where $X_i\neq X_j$ for all $i\neq j$, there exists $F\in \mathcal{T}_{\mathcal{G}, \mathcal{H}}$ such that $\|F(X_i)-Y_i\|<\varepsilon$ for all $i$.
    \item \textbf{Approximation of Tensor-Type Functions}:
    For any continuous, increasing function $h:\mathbb{R}\to\mathbb{R}$ and any $\varepsilon > 0$, there exists $F\in\mathcal{T}_{\mathcal{G}, \mathcal{H}}$ such that
    \begin{equation}
        \|F-h^{(d\times n)}\|_{C(K)}\leq \varepsilon,
    \end{equation}
    where $h^{(d\times n)}$ is the coordinate-wise extension of $h$ to $\mathbb{R}^{d\times n}$, given by
    \begin{equation}
        (h^{(d\times n)}(X))_{ij}=h(X_{ij}), \quad 1\leq i\leq d, \ 1\leq j\leq n.
    \end{equation}
\end{itemize}
The intuition behind this framework is that the interpolation property ensures pointwise matching of the target function, while the approximation of tensor-type functions controls the error outside the interpolated samples.

For the function family $\tilde{\mathcal{H}}$ mapping $\mathbb{R}^d\to\mathbb{R}^d$, define
\begin{equation}
    \mathcal{A}_{\tilde{\mathcal{H}}}:=\{(\operatorname{Id}+h_m)\circ \cdots (\operatorname{Id}+h_1) \mid m\in\mathbb{N}_+, \ h_i\in \tilde{\mathcal{H}} \text{ for all } i\}.
\end{equation}
Results from~\cite{cheng2025interpolation} establish that $\mathcal{A}_{\tilde{\mathcal{H}}}$ satisfies the universal interpolation property (UIP) for functions $\mathbb{R}^d\to\mathbb{R}^d$ if $\tilde{\mathcal{H}}$ meets the conditions in \Cref{thm:main}. This implies that $\mathcal{T}_{\mathcal{G}, \mathcal{H}}$ can interpolate any dataset $\{(X_i, Y_i)\}_{i=1}^N$ if all tokens in each $X_i$ are distinct, simply by setting the attention block to zero. Moreover, the approximation of tensor-type functions follows in the same way. However, token-wise feature maps alone cannot interpolate data containing duplicate tokens.

Thus, attention mechanisms play a crucial role by generating contextual information to distinguish tokens before token-wise feature maps perform interpolation. Using results from~\cite{cheng2025interpolation}, proving the UAP of $\mathcal{T}_{\mathcal{G}, \mathcal{H}}$ reduces to verifying the distinguishability of tokens across finite data samples. For $G$-UAP, since the target functions are permutation equivariant, token distinguishability is also be analyzed under permutation equivariance. See~\Cref{sec:proof_main} in the appendix for a detailed proof of \Cref{thm:main}.

For \Cref{prop:analytic}, assume that Condition 2 in \Cref{thm:main} fails. Then, there exists $N$ samples $\{X_i\}_{i=1}^N$ from different orbits under the $G$-action, but for any $m$ and $g\in (\operatorname{Id}+\mathcal{G})^m$, there exist indices $i,j\in [N]$ such that at least one token in $g(X_i)$ is identical to a token in $g(X_j)$. This can be written as
\begin{equation}
    \label{eq:fail}
    \Pi_{i,j}:=\prod_{l_1,l_2} \|[g(X_i)]_{l_1}-[g(X_j)]_{l_2}\|_2^2=0.
\end{equation}
Since $\mathcal{G}$ is parametric and analytic in $\theta$, the function $\Pi_{i,j}$ is also analytic in $\theta$. As the zero set of a nontrivial analytic function has measure zero, this implies that for some $i,j$, $\Pi_{i,j}$ is identically zero, meaning that $\mathcal{G}$ fails to distinguish tokens in $X_i$ and $X_j$. This argument reduces Condition 2 in \Cref{thm:main} to the case $N=2$. The existence of a uniform $m$ follows from the same reasoning.

% \subsection{Proof Idea for \Cref{thm:kernel} and \Cref{prop:sparse}}

% \ql{I think this part is optional, or at least move it behind the result}

% \Cref{thm:kernel} and Cref{prop:sparse} rely on the following technical lemma:

% \begin{lemma}
%     Let $r$ be a kernel function satisfying the conditions in \Cref{thm:kernel}.
%     Suppose $\{a_1, \dots, a_r\}\subset \mathbb{R}^d\setminus \{0\}$ and $\{b_1,\dots, b_s\}\subset \mathbb{R}^d\setminus \{0\}$ are two sequences of distinct tokens. Suppose that for some indices $r^\prime\in [r]$ and $s^\prime\in [s]$, the following equality holds for all choices of $W_Q, W_K, W_V$:
%     \begin{equation}
%         % \label{eq:lem_equal_atten}
%         a_{r^\prime}+\sum_{j=1}^r\left(\frac{ k(W_Q a_{r^\prime}, W_K a_j)}{\sum_{l=1}^r k(W_Q a_{r^\prime}, W_K a_l)} W_V a_j\right) =
%         b_{s^\prime}+\sum_{j=1}^s\left(\frac{ k(W_Q b_{s^\prime}, W_K b_j)}{\sum_{l=1}^s k(W_Q b_{s^\prime}, W_K b_l)} W_V b_j\right).
%     \end{equation}
%     Then, $r=s$, and there exists a permutation $\sigma\in{S}_n$ with $\sigma(r^\prime)=s^\prime$ such that $a_i=b_{\sigma(i)}$ for all $i$.
% \end{lemma}

% \Cref{lem:equal_atten} states that if a kernel function satisfies the conditions in Theorem~\ref{thm:kernel}, then the attention mechanism generates distinct contextual embeddings for distinct token sets.

% For the sparse transformer, \Cref{lem:equal_atten} indicates that if the $i_1$-the token of $X$ and $i_2$-th token of $i_2$ cannot be distinguished by $\operatorname{Id}+\mathcal G$, we can deduce that the set of tokens with indices in $\mathcal N(i_1)$ of $X$ and those in $\mathcal N(i_2)$ of $Y$ must be the same. Then, we can choose a common token in the two sets, and repeat the same discussion for $(\operatorname{Id}+\mathcal G)^2$. Since $\Gamma$ is strongly connected, this process can be repeated to cover all the indices in $[n]$, resulting in a permutation $\sigma$, such that $\sigma$ is an automorphism of $\Gamma$ and $[X]_i=[Y]_{\sigma(i)}$ for all $i$, i.e. $X$ and $Y$ are in the same orbit of $\operatorname{Aut}(\Gamma)$.
% See section~\ref{} in the appendix for the detailed proof.
\end{comment}


\bibliographystyle{plain}
\bibliography{ref}
\clearpage


%%%%%%%%%%%%%%%%%%%%%%%%%%%%%%%%%%%%%%%%%%%%%%%%%%%%%%%%%%%%

\appendix
\section{Detailed related works}
\label{sec:related_work}
% \section{Approximation results for ResNets}
% As a fundamental aspect in deep learning theory, the expressive capability of deep neural networks has been extensively studied. The universal approximation results for feedforward ~\cite{hornik1991approximation}

\paragraph{Approximation results for transformers}
Since the introduction of the transformer architecture in~\cite{vaswani2017attention}, numerous studies have investigated its approximation properties. The universal approximation property (UAP) of the original transformer with softmax attention as a fixed length sequence-to-sequence model was established in~\cite{yun2019transformers}. This constructive approach was later extended to transformers with certain sparse attention mechanisms~\cite{yun2020n, zaheer2020big}. In~\cite{kajitsukatransformers}, the authors proposed a new construction demonstrating that transformers with a single attention layer and sufficiently deep feedforward networks can achieve UAP. Similarly,\cite{alberti2023sumformer} showed that with increased hidden dimensions, two variants of transformers, i.e. LinFormer~\cite{wang2020linformer} and Performer~\cite{choromanski2020rethinking}, satisfy the UAP.
In contrast to these constructive methods for different architectures, our results offer a unified framework for establishing the UAP of various transformer models without relying on explicit constructions. 
Beyond these works treating transformers as a fixed length sequence-to-sequence model, there are also studies handling transformers with variable-length inputs by considering the input sequence as an empirical measure~\cite{furuya2024transformers,geshkovski2024measure,geshkovski2025mathematical}. The universal interpolation and universal approximation properties under this viewpoint were established with proper assumptions. Compared to these results, we still consider transformers as sequence-to-sequence models in this work to offer a direct analysis for different transformer architectures. Recently, there are other works studying the UAP of attention-only architectures~\cite{liu2025attention, hu2025universal}, indicating that softmax attention alone can also achive strong approximation power.
Other studies have explored the UAP of transformers under alternative settings, such as in-context learning, prompting, and constrained scenarios~\cite{kratsios2021universal, luo2022your, furuya2024transformers, petrov2024prompting,hu2025fundamental}. Another researchline is the Turing completeness of transformers~\cite{perez2021attention, wei2022statistically}.
Besides these UAP results, there are also works providing the approximation rates of transformers. For instance,\cite{jiang2024approximation} provides explicit rates over a dense subset of sequence-to-sequence functions; \cite{wang2024understanding} derives rates for target functions with structured memory; and~\cite{takakura2023approximation} characterizes the approximation rate in terms of function smoothness for transformers with infinitely long inputs.

% \paragraph{Approximation results for transformers}
% Since the transformer architecture was proposed in~\cite{vaswani2017attention}, a series of works have studied its approximation properties. 
% The universal approximation property (UAP) for the original transformer with softmax attention was established in~\cite{yun2019transformers}. 
% The constrcted approach was attended to transformers with certain sparse-types~\cite{yun2020n,zaheer2020big}. In~\cite{kajitsukatransformers}, the authors used a new construction to show that transformers with a single attention layer and enough feedforward layers can achieve UAP. In~\cite{alberti2023sumformer}, the author shows that with additional hidden dimensions, LinFormer~\cite{wang2020linformer} and Performer~\cite{choromanski2020rethinking} can achieve UAP. 
% Compared to these constructive approaches on different types of transformers, our results provide a unified framework for establishing the UAP of different transformer architectures without the need for explicit constructions. 
% In~\cite{kratsios2021universal, luo2022your, furuya2024transformers, petrov2024prompting}, the UAP of transformers are studied in different settings, including in-context learning, prompting and under certain constraints. There is another research line that studies the Turing completeness of transformers~\cite{perez2021attention,perez2021attention,wei2022statistically}. Besides the UAP results, there are also works that give approximation rate results of transformers. For example,~\cite{jiang2024approximation} provides explicit approximation rates in a dense subset of sequence-to-sequence funtions. \cite{wang2024understanding} provides approximation rates for target functions with certain memory structures. \cite{takakura2023approximation} provides rate characterized by the smoothness for transformers with infinite length inputs.

\paragraph{Approximation under symmetry}
The study of approximation under functional symmetries has been explored in various works. 
In~\cite{yarotsky2022universal}, the universal approximation of functions invariant under compact group or translation actions was analyzed using shallow neural networks. 
In~\cite{cohen2016group,finzi2020generalizing}, convolutional structures were proposed to approximate equivariant functions.
In~\cite{ravanbakhsh2020universal,maron2019universality,li2024deep}, the universal approximation under symmetry  using deep neural networks was investigated.
Notably,~\cite{li2024deep} provides a general sufficient condition for the action of any transitive subgroup of $S_n$ on coordinates (1-dimensional tokens). In contrast, our work addresses approximation under symmetry in a broader setting, considering group actions on $d$-dimensional tokens rather than coordinates, thereby extending the analysis in \cite{li2024deep} to non-transitive permutation groups. In~\cite{agrachev2025generic}, the authors provide a general framework on the ensemble controllability of control systems under symmetry. They also show that systems that can interpolate arbitrarily many samples under symmetry are generic in certain topology. Compare to their genericity results, our results provide a verifiable sufficient condition for UAP of specific architectures, allowing direct applications to various transformer architectures.

\paragraph{Transformer variants}
Beyond the original Transformer, numerous architectural variants have been developed to improve efficiency, scalability, or adaptability. These include sparse attention mechanisms~\cite{beltagy2020longformer,correia2019adaptively,kitaev2020reformer,zaheer2020big,guo2019star}, low-rank and kernel-based approximations of attention~\cite{wang2020linformer,choromanski2020rethinking,xiong2021nystromformer}, as well as other architectural modifications~\cite{liu2023blockwise,haldar2024baku,mamega,zhai2021attention}.
Another related line of work explores parameter-efficient fine-tuning methods for large Transformer models~\cite{hu2022lora,liu2024dora,meng2024pissa,jia2022visual,zhang2024parameter,zhangweight}, which aim to adapt pretrained networks to downstream tasks with minimal additional parameters.
In this paper, we establish a general sufficient condition for the universal approximation property (UAP) of various Transformer variants. As demonstrated in \Cref{sec:app}, our framework can be readily verified for many existing architectures, and potentially extended to other designs not included as well.


% IO-aware exact kernels (FlashAttention~\cite{dao2022flashattention,dao2023flashattention2}); positional schemes for long-range generalization (Transformer-XL/T5~\cite{dai2019transformer,raffel2020exploring}; RoPE/ALiBi~\cite{su2021rope,press2021alibi}); mixture-of-experts scaling (GShard/Switch~\cite{lepikhin2021gshard,fedus2021switch}); head/key–value sharing (multi-query/grouped-query~\cite{shazeer2019fast,ainslie2023gqa}); feed-forward/normalization refinements (GEGLU/SwiGLU; Pre-LN/RMSNorm~\cite{shazeer2020glu,xiong2020layernorm,zhang2019rmsnorm}); hierarchical/local–global hybrids in vision (ViT/Swin~\cite{dosovitskiy2021vit,liu2021swin}); and retrieval-/memory-augmented models (RAG/RETRO~\cite{lewis2020rag,borgeaud2022retro}). These directions are complementary to our analysis, which formulates UAP criteria at fixed context length without committing to a specific attention instantiation.



\section{Proof of \Cref{thm:main} and \Cref{prop:analytic}}

% \subsection{Proof of Theorem~\ref{thm:main} and Proposition~\ref{prop:analytic}}
\subsection{Proof of \Cref{thm:main}}
\label{sec:proof_main}
In the following, $\|\cdot\|_2$ denotes the $\ell^2$-norm, for both  vectors in $\mathbb R^{d\times n}$ or  $\mathbb R^{d}$. We begin by proving the following interpolation property of $\mathcal{T}_{\mathcal{G}, \mathcal{H}}$:

% \lt{What is $\|$ for a point?}
\begin{proposition}[\textbf{Interpolation Property}]
    \label{prop:interpolation}
 Suppose $\mathcal G$ and $\mathcal H$ satisfy the condition in \Cref{thm:main} for group $G$.
 Consider any $G$-equivariant continuous function $F:\mathbb R^{d\times n}\to\mathbb R^{d\times n}$.
 Then, for any $\varepsilon >0$ and $\{X_i\}_{i=1}^N\subset \mathbb R^{d\times n}$, there exists $\hat F\in \mathcal T_{\mathcal G, \mathcal H}$ such that:
 \begin{itemize}
    \item  $\|\hat F(X_i)-F(X_i)\|_2<\varepsilon$, if $X_i$ is in general position(defined in~\Cref{sec:preliminaries}). 
    \item $\displaystyle\|\hat F(X_i)\|_2< n\cdot \max_{i}\{\|F(X_i)\|_2\} + 2 \varepsilon$, if $X_i$ is not in general position.
 \end{itemize}
%  Consider any finite set $\{(X_i, Y_i)\}_{i=1}^N\subset K\times \mathbb{R}^{d\times n}$
%     data where $X_i$ are all in general positions, and
%     Then, for any $\varepsilon>0$ and any finite dataset $\{(X_i, Y_i)\}_{i=1}^N\subset K\times \mathbb{R}^{d\times n}$. where $X_i\neq X_j$ for all $i\neq j$, there exists $F\in \mathcal{T}_{\mathcal{G}, \mathcal{H}}$ such that $\|F(X_i)-Y_i\|<\varepsilon$ for all $i$.
\end{proposition}

\begin{proof}
                                     Since $F$ and functions in $\mathcal T_{\mathcal G, \mathcal H}$ are $G$-equivariant, we only need to consider the case when $X_i$ are from distinct orbits under the $G$-action.
    Moreover, we can assume that $X_i$ are in general position for $i=1,\cdots, M$, and $X_i$ are not in general position for $i=M+1, \cdots, N$.

    By the token distinguishability condition, there exist $m$ and $g\in (\operatorname{Id}+\mathcal{G})^m$ such that the tokens of $g(X_i)$ are all distinct for $i=1,\cdots, M$.
    We denote $x_1,\cdots, x_{Mn}$ as the distinct tokens in $g(X_1),\cdots, g(X_M)$.
    For each $j\le Mn$, suppose $x_j=[X_l]_k$ for some $(l,k)\in [n]\times [M]$, we denote $y_j$ as its corresponding token $[F(X_l)]_k$. Moreover, we denote
 $x_{Mn+1},\cdots, x_{Mn+J}\in\mathbb R^d$ as the distinct tokens in $g(X_{M+1}),\cdots, g(X_N)$ that are different from $x_1,\cdots, x_{Mn}$. For each $j>Mn$, we denote $y_j=0\in\mathbb R^d$. Then, we get a set of $d$-dimensional pairs $\{(x_j, y_j)\}_{j=1}^{Mn+J}\subset \mathbb R^{d}\times \mathbb R^d$ where all $x_j$ are all distinct. Since $\mathcal H$ satisfies the non-linear affine invariance condition, according to the main result in~\cite{cheng2025interpolation}, we know that there exists a function
 \begin{equation}
    f\in \{(\mathrm{Id}+f_k)\circ \cdots \circ (\mathrm{Id}+f_1)\mid k\in \mathbb N_+,  f_1,\cdots, f_k\in\mathbb \mathcal H\}
 \end{equation}
such that
\begin{equation}
    \|f(x_j)-y_j\|_2\le \frac{1}{n}\varepsilon, \quad j=1,\cdots, Mn+J.
\end{equation}
Denote $f^{\otimes n}$ as the token-wise extension of $f$ to $\mathbb R^{d\times n}$. For each $i\le M$, we have
\begin{equation}
    \begin{aligned}
    \|f^{\otimes n}(g(X_i))-F(X_i)\|_2&=\left\|\Big(\big(f(g[X_i]_1)-[F(X_i)]_1\big), \cdots, \big(f(g[X_i]_n)-[F(X_i)]_n\big)\Big)\right\|_2\\
    &\le n\cdot \max_{j}\|f(x_j)-y_j\|_2\le \varepsilon.
    \end{aligned}
\end{equation}
For each $i\ge M+1$, we have
\begin{equation}
    \begin{aligned}
        \|f^{\otimes n}(g(X_i))\|&=\left\|\Big(f(g[X_i]_1), \cdots, f(g[X_i]_n)\Big)\right\|_2\\
        &\le \sum_{l: g[X_i]_l\in \{x_1,\cdots, x_{Mn}\}}\|f(g[X_i]_l)\|_2 + \sum_{l: g[X_i]_l\in \{x_{Mn+1},\cdots, x_{Mn+J}\}}\|f(g[X_i]_l)\|_2\\
        &\le n\cdot \max_{j}\|f(x_j)\|_2 + n\cdot \max_{j>Mn}\|f(x_j)\|_2\\
        & \le n\cdot \left(\max_{i}\{\|F(X_i)\|_2\}+\frac{\varepsilon}{n} \right) + \varepsilon = n\cdot \max_{i}\{\|F(X_i)\|_2\} + 2\varepsilon.
        \end{aligned}
\end{equation}
Therefore, $\hat F=f^{\otimes n}\circ g \in \mathcal T_{\mathcal G, \mathcal H}$ satisfies the interpolation property.
\end{proof}

\begin{proposition}[\textbf{Approximation of Tensor-Type Functions}]
    \label{prop:tensor_type_approx}
    Suppose $\mathcal G$ and $\mathcal H$ satisfy the condition in \Cref{thm:main} for group $G$.
    Consider any continuous, increasing function $h:\mathbb{R}\to\mathbb{R}$ and any $\varepsilon > 0$.
    Then, there exists $F\in\mathcal{T}_{\mathcal{G}, \mathcal{H}}$ such that
    \begin{equation}
        \|F-h^{(d\times n)}\|_{C(K)}\leq \varepsilon,
    \end{equation}
    where $h^{(d\times n)}$ is the coordinate-wise extension of $h$ to $\mathbb{R}^{d\times n}$, given by
    \begin{equation}
        (h^{(d\times n)}(X))_{ij}=h(X_{ij}), \quad 1\leq i\leq d, \ 1\leq j\leq n.
    \end{equation}
\end{proposition}
\begin{proof}
    \Cref{prop:tensor_type_approx} directly follows from the proof of Proposition 4.11 ~\cite{li2022deep} and Theorem 2.6 in~\cite{cheng2025interpolation}.
\end{proof}

\begin{proposition}[Corollary of Main results of ~\cite{cheng2025interpolation}]
    \label{prop:interpolation_2}
    Let $d\ge 2$ and $\mathcal H$ be a family of maps from $\mathbb R^d$ to $\mathbb R^d$ that satisfies the non-linearity and affine-invariance condition in \Cref{thm:main}. Then, for any $  \{(x_i,y_i)\}_{i=1}^N\subset \mathbb R^d\times \mathbb R^d$ with $x_i\neq x_j$ for all $i\neq j$ and $\varepsilon>0$, there exists
    \begin{equation}
    f\in \{(\mathrm{Id}+f_k)\circ \cdots \circ (\mathrm{Id}+f_1)\mid k\in \mathbb N_+,  f_1,\cdots, f_k\in \mathcal H \}
    \end{equation}
    such that
    \begin{equation}
        \|f(x_i)-y_i\|\le \varepsilon, \quad i=1,\cdots, N.
    \end{equation}
\end{proposition}
\begin{proof}
    ~\Cref{prop:interpolation_2} is a direct corollary of Theorem 2.6 in~\cite{cheng2025interpolation}.
\end{proof}

\begin{proof}[Proof of \Cref{thm:main}]
    The approach of the proof is similar to the main theorem in~\cite{li2022deep} and ~\cite{li2024deep}.

    We assume without loss of generality that
    $
    K = [-s,s]^{d\times n}
    $
    is a hypercube in $\mathbb{R}^{d\times n}$. Our target is to show that for any $\varepsilon >0$, we can find function $\hat F\in \mathcal{T}_{\mathcal{G},\mathcal{H}}$ such that
    $
    \|\hat f-F\|_{L^p(K)}\le \varepsilon
    $.

    \medskip
    \textbf{Step 1.} For each multi-index $\mathbf i=(i_{kl})_{k\in [d], l\in [n]}\in \mathbb Z^{d\times n}$ and $\delta>0$, we define the grid cells:
    \begin{equation}
        \square_{\mathbf i, \delta}:=\left\{X\in \mathbb R^{d\times n}\mid X_{kl}\in [i_{kl}\delta, (i_{kl}+1)\delta], \text{ for all } k\in [d], l\in [n]\right\}.
    \end{equation}
    We denote $\bm p_{\mathbf i,\delta}:=\delta\mathbf i$ as a corner point of $\square_{\mathbf i, \delta}$, and $\chi_{\bm i, \delta}$ as the characteristic function of $\square_{\bm i, \delta}$.
    Since $F$ is continuous, there exists $\delta> 0$ such that it has a piece-wise constant approximation
    \begin{equation}
        \tilde F:=\sum_{\mathbf i} F(\bm p_{\mathbf i})\chi_{\mathbf i, \delta},
    \end{equation}
    such that
    \begin{equation}
        \|F-\tilde F\|_{L^p(K)}\le \frac{\varepsilon}{4}.
    \end{equation}

    \medskip
    \noindent\textbf{Step 2.}
    Apply \Cref{prop:interpolation_2} to the set of grid points
    \[
    \Big\{ \bm{p}_{\mathbf{i},\delta} \mid \mathbf{i}_{kl}\in\Big\{\lceil -s/\delta \rceil, \dots, \lfloor s/\delta \rfloor\Big\} \Big\}.
    \]
    Then, for any $\gamma > 0$, there exists a function $\bar{F} \in \mathcal{T}_{\mathcal{G},\mathcal{H}}$ such that:
    \begin{itemize}
        \item If $\bm{p}_{\mathbf{i},\delta}$ is in general position, then
        \[
        \|\bar{F}(\bm{p}_{\mathbf{i},\delta}) - F(\bm{p}_{\mathbf{i},\delta})\| < \gamma.
        \]
        \item Otherwise,
        \[
        \|\bar{F}(\bm{p}_{\mathbf{i},\delta})\| < n\cdot\max_i\{\|F(\bm{p}_{\mathbf{i},\delta})\|\} + 2\gamma\le n\|F\|_{C(K)} + 2\gamma.
        \]
    \end{itemize}


    \medskip
    \noindent \textbf{Step 3.}
    For any $\alpha\in (0,1)$, define
    \begin{equation}
        \square_{\mathbf i, \delta}^\alpha:=\left\{X\in \mathbb R^{d\times n}\mid X_{kl}\in [i_{kl}\delta, i_{kl}\delta+\alpha \delta], \text{ for all } k\in [d], l\in [n]\right\},
    \end{equation}
    as the shrunk hypercube of $\square_{\mathbf i, \delta}$ with side length $\alpha\delta$.
    We consider the map $h_{\alpha,\delta}:\mathbb R\to \mathbb R$ defined as:
    \begin{equation}
        h_{\alpha,\delta}(x)=
        \begin{cases}
        i\delta, & \text{if } x\in [i\delta,\, i\delta+\alpha\delta],\quad i\in\mathbb{Z},\\[1mm]
        i\delta+\dfrac{x-i\delta-\alpha\delta}{1-\alpha}, & \text{if } x\in [i\delta+\alpha\delta,\, (i+1)\delta],\quad i\in\mathbb{Z}.
        \end{cases}
    \end{equation}
    On the interval \([i\delta,\, i\delta+\alpha\delta]\), the function \(h_{\alpha,\delta}(x)\) remains constant at \(i\delta\). Then, as \(x\) increases from \(i\delta+\alpha\delta\) to \((i+1)\delta\), the function increases linearly from \(i\delta\) to \((i+1)\delta\). Notice that $h_{\alpha,\delta}^{(d\times n)}$ is continuous and has constant value $\bm p_{\mathbf i}$ on $\square_{\mathbf i, \delta}^\alpha$ for each $\mathbf i$.
    Since \(h_{\alpha,\delta}\) is continuous and increasing, by \Cref{prop:tensor_type_approx}, for any $\rho>0$,  there exists a function \(H_{\alpha,\delta}\in\mathcal{T}_{\mathcal{G}, \mathcal{H}}\) such that
    \begin{equation}
        \|H_{\alpha,\delta}-h_{\alpha,\delta}^{(d\times n)}\|_{C(K)}< \rho.
    \end{equation}

    \medskip
    \textbf{Step 4.}
    Now we can estimate the error of the composition $\|\bar{F}\circ H_{\alpha,\delta}-F\|_{L^p(K)}$. We define
    \[
    K^\alpha = \Big(\bigcup_{\mathbf{i}} \square_{\mathbf{i},\delta}^\alpha\Big) \cap K,
    \]
and denote $K_1^\alpha$ as the union of the grid cells in $K^\alpha$ whose corner point $\bm p_{\mathbf i}$ are in general position, and $K_2^\alpha$ as the union of the grid cells in $K^\alpha$ whose $\bm p_{\mathbf i}$ are not in general position. Then, we have
    \begin{enumerate}
        \item On $K_1^\alpha$, we have
        \begin{equation}
        \|\bar{F}\circ h_{\alpha,\delta}^{(d\times n)} - \tilde{F}\|_{L^p(K_1^\alpha)} < \gamma (m(K_1^\alpha))^{\frac{1}{p}}.
        \end{equation}
        \item For $K^\alpha_2$,  the number of $\mathbf i$ with $\bm p_{\mathbf i, \delta}$ not in general position is at most $\frac{n(n-1)}{2}(2s/\delta +1)^{(n-1)d}$, we have the measure of $K_2^\alpha = \mathcal O(\delta^{d})$.
        Therefore,
        \begin{equation}
        \|\bar{F}\circ h_{\alpha,\delta}^{(d\times n)} - \tilde{F}\|_{L^p(K_2^\alpha)} < (n\cdot \|F\|_{C(K)} + 2\gamma) \mathcal O(\delta^d).
        \end{equation}

    \end{enumerate}
    Therefore, we can choose $\delta $ and $\gamma$ sufficiently small such that
    \begin{equation}
            \|\bar{F}\circ h_{\alpha,\delta}^{(d\times n)} - \tilde{F}\|_{L^p(K^\alpha)} < \frac{\varepsilon}{4}.
    \end{equation}

On $K\setminus K^\alpha$, by choosing $\alpha$ sufficiently close to $1$, we make $m(K\setminus K^\alpha)$ arbitrarily small. Since $\bar F\circ h_{\alpha,\delta}^{(d\times n)}$ and $\tilde F$ is bounded on $K$, the following can be guaranteed:
    \begin{equation}
    \|\bar{F}\circ h_{\alpha,\delta}^{(d\times n)} - \tilde{F}\|_{L^p(K\setminus K^\alpha)} < \frac{\varepsilon}{4}.
    \end{equation}

    Since $\bar F$ is uniformly continuous on $K$, there exists $\rho>0$ such that for any $X, Y\in K$ with  $\|X-Y\|<\rho$, we have $\|\bar F(X)-\bar F(Y)\|<\kappa:=\varepsilon/(4(m(K))^{\frac{1}{p}})$. Therefore, after determining $\alpha, \delta$ and $\gamma$, we can choose $\rho$ such that
     \begin{equation}
        \|\bar F\circ H_{\alpha,\delta}-F\|_{L^p(K)}\le (m(K))^{\frac{1}{p}}\kappa \le \frac{\varepsilon}{4}.
     \end{equation}After determining $\delta$ and $\alpha$, by step 3, we can choose $\rho$ sufficiently small such that 
     \begin{equation}
        \| \bar{F}\circ H_{\alpha,\delta} - \bar{F}\circ h_{\alpha,\delta}^{(d\times n)}\|_{L^p(K)}\le \rho \operatorname{\operatorname{Lip}(\bar F)}<\frac{\varepsilon}{4}.
     \end{equation}

    Recall also that
    \[
    \|\tilde{F} - F\|_{L^p(K)} < \frac{\varepsilon}{4}.
    \]
    Thus, by the triangle inequality,
    \begin{equation}
    \begin{aligned}
    \| \bar{F}\circ H_{\alpha,\delta} - F\|_{L^p(K)}& \le \| \bar{F}\circ H_{\alpha,\delta} - \bar{F}\circ h_{\alpha,\delta}^{(d\times n)}\|_{L^p(K)}+\|\bar{F}\circ h_{\alpha,\delta}^{(d\times n)}-\tilde F \|_{L^p(K^\alpha)} \\
    & + \|\bar{F}\circ h_{\alpha,\delta}^{(d\times n)}-\tilde F \|_{L^p(K\setminus K^\alpha)} + \|\tilde F - F\|_{L^p(K)}\\
    & < \frac{\varepsilon}{4} + \frac{\varepsilon}{4} + \frac{\varepsilon}{4} + \frac{\varepsilon}{4} = \varepsilon,
    \end{aligned}
    \end{equation}
    which completes the proof.
    \end{proof}

% \begin{proof}[\textbf{Proof of Theorem~\ref{thm:main}}]
%     Without loss of generality, we assume that $K=[-s,s]^{d\times n}$ is a hypercube in $\mathbb R^{d\times n}$. For each multi-index $\mathbf i=(i_{kl})_{k\in [d], l\in [n]}\in \mathbb Z^{d\times n}$ and $\delta>0$, we denote
%     \begin{equation}
%         \square_{\mathbf i, \delta}:=\left\{X\in \mathbb R^{d\times n}\mid X_{kl}\in [i_{kl}\delta, (i_{kl}+1)\delta], \text{ for all } k\in [d], l\in [n]\right\}.
%     \end{equation}
%     We denote $\bm p_{\mathbf i,\delta}:=\delta\mathbf i$ as a corner point of $\square_{\mathbf i, \delta}$, and $\chi_{\bm i, \delta}$ as the characteristic function of $\square_{\bm i, \delta}$.
%     Since $F$ is continuous, there exists $\delta\ge 0$ such that it has a piece-wise constant approximation
%     \begin{equation}
%         \tilde F:=\sum_{\mathbf i} F(\bm p_{\mathbf i})\chi_{\mathbf i, \delta},
%     \end{equation}
%     such that
%     \begin{equation}
%         \|F-\tilde F\|_{L^p(K)}\le \frac{\varepsilon}{4}.
%     \end{equation}
%     On the other hand, for any $\alpha\in (0,1)$, denote
%     \begin{equation}
%         \square_{\mathbf i, \delta}^\alpha:=\left\{X\in \mathbb R^{d\times n}\mid X_{kl}\in [i_{kl}\delta, i_{kl}\delta+\alpha \delta], \text{ for all } k\in [d], l\in [n]\right\},
%     \end{equation}
%     as the shrunk hypercube of $\square_{\mathbf i, \delta}$ with side length $\alpha\delta$.
%     We consider the map $h_{\alpha,\delta}:\mathbb R\to \mathbb R$ defined as:
%     \begin{equation}
%         h_{\alpha,\delta}(x)=
%         \begin{cases}
%         i\delta, & \text{if } x\in [i\delta,\, i\delta+\alpha\delta],\quad i\in\mathbb{Z},\\[1mm]
%         i\delta+\dfrac{x-i\delta-\alpha\delta}{1-\alpha}, & \text{if } x\in [i\delta+\alpha\delta,\, (i+1)\delta],\quad i\in\mathbb{Z}.
%         \end{cases}
%     \end{equation}
%     On the interval \([i\delta,\, i\delta+\alpha\delta]\), the function \(h_{\alpha,\delta}(x)\) remains constant at \(i\delta\). Then, as \(x\) increases from \(i\delta+\alpha\delta\) to \((i+1)\delta\), the function increases linearly from \(i\delta\) to \((i+1)\delta\). Notice that $h_{\alpha,\delta}^{(d\times n)}$ is continuous and has constant value $\bm p_{\mathbf i}$ on $\square_{\mathbf i, \delta}^\alpha$ for each $\mathbf i$.
%     Since \(h_{\alpha,\delta}\) is continuous and increasing, by Proposition~\ref{prop:tensor_type_approx}, for any $\rho>0$,  there exists a function \(H_{\alpha,\delta}\in\mathcal{T}_{\mathcal{G}, \mathcal{H}}\) such that
%     \begin{equation}
%         \|H_{\alpha,\delta}-h_{\alpha,\delta}^{(d\times n)}\|_{C(K)}< \rho.
%     \end{equation}
%     Consider applying the interpolation property in Proposition~\ref{prop:interpolation} to the set
%     \begin{equation}
%         \left\{ \bm p_{\mathbf i, \delta}\;\middle|\; \mathbf i_{kl} \in \left\{ \lceil -\frac{s}{\delta} \rceil, \lceil -\frac{s}{\delta} \rceil + 1, \dots, \lfloor \frac{s}{\delta} \rfloor \right\}, \, \forall\, k\in [d],\, l\in [n] \right\},
%     \end{equation} we know that for any $\gamma>0$ there exists $\bar F \in \mathcal T_{\mathcal G, \mathcal H}$ such that
%     \begin{itemize}
%         \item  $\|\bar F(\bm p_{\mathbf i, \delta})-F(\bm p_{\mathbf i, \delta})\|<\gamma$, if $\bm p_{\mathbf i, \delta}\in\mathbb R^{d\times n}$ is in general position.
%         \item $\displaystyle\|\bar F(\bm p_{\mathbf i, \delta})\|< n\cdot \max_{i}\{\|F(\bm p_{\mathbf i, \delta})\|\} + 2 \gamma$, if $\bm p_{\mathbf i, \delta}$ is not in general position.
%      \end{itemize}
%      Then, considering the function $\bar F\circ H_{\alpha,\delta}$, we have
%      \begin{equation}
%         \label{eq:estimate}
%                 \|\bar F\circ H_{\alpha,\delta}-F\|_{L^p(K)}\le \|\bar F\circ H_{\alpha,\delta}-\bar F\circ h_{\alpha,\delta}^{(d\times n)}\|_{L^p(K)}+\|\bar F\circ h_{\alpha,\delta}^{(d\times n)}-F\|_{L^p(K)}
%      \end{equation}
%      For the first term, since $\bar F$ is uniformly continuous on $K$, there exists $\rho>0$ such that for any $X, Y\in K$ with  $\|X-Y\|<\rho$, we have $\|\bar F(X)-\bar F(Y)\|<\bar\gamma$. Therefore, we can choose such $\rho$ and we have
%      \begin{equation}
%         \|\bar F\circ H_{\alpha,\delta}-F\|_{L^p(K)}\le (m(K))^{\frac{1}{p}}\bar \gamma \le \frac{\varepsilon}{4}.
%      \end{equation}
%      For the second term in~~\eqref{eq:estimate}, we have
%      \begin{equation}
%         \label{eq:estimate2}
%         \begin{aligned}
%         \|\bar F\circ h_{\alpha,\delta}^{(d\times n)}-F\|_{L^p(K)}&\le \|\bar F\circ h_{\alpha,\delta}^{(d\times n)}-\tilde F\|_{L^p(K)}+\|\tilde F-F\|_{L^p(K)}\\
%         &\le \|\bar F\circ h_{\alpha,\delta}^{(d\times n)}-\tilde F\|_{L^p(K^\alpha)}+\|\bar F\circ h_{\alpha,\delta}^{(d\times n)}-\tilde F\|_{L^p(K\setminus K^\alpha)}+\frac{\varepsilon}{4},
%         \end{aligned}
%      \end{equation}
%      where $K^\alpha=\left(\bigcup_{\mathbf i} \square_{\mathbf i, \delta}^\alpha \right)\cap K$.
%     Since both $\bar F$ and $\tilde F$ are bounded on $K$, We can choose $\alpha$ close to $1$ enough such that the measure of $K\setminus K^\alpha$ is small enough, and that
%     \begin{equation}
%         \|\bar F\circ h_{\alpha,\delta}^{(d\times n)}-\tilde F\|_{L^p(K\setminus K^\alpha)}\le \frac{\varepsilon}{4}.
%     \end{equation}
%     For the first term in~~\eqref{eq:estimate2}, we have
%     \begin{equation}
%         \| F\circ h_{\alpha,\delta}^{(d\times n)}(X)-\tilde F(X)\|=\|\bar F(\bm p_{\mathbf i, \delta})-F(\bm p_{\mathbf i, \delta})\|\le \begin{cases}
%             \gamma, & \text{if } \bm p_{\mathbf i, \delta} \text{ is in general position},\\
%             n\cdot \max_{i}\{\|F(\bm p_{\mathbf i, \delta})\| \}+2\gamma, & \text{otherwise}.
%         \end{cases}
%     \end{equation}
%     on $\square_{\mathbf i, \delta}^\alpha$. We can calculate that the number of $\mathbf i$ with $\bm p_{\mathbf i, \delta}$ not in general position is at most $\frac{n(n-1)}{2}(2s/\delta +1)^{(n-1)d}$, we have the measure of these shrunk cubes is in order $\mathcal O(\delta^{d})$. Therefore, we can choose $\delta$ and $\gamma $ small enough such that
%     \begin{equation}
%         \|\bar F\circ h_{\alpha,\delta}^{(d\times n)}-\tilde F\|_{L^p(K^\alpha)}\le \frac{\varepsilon}{4}.
%     \end{equation}
%     Combine all the estimates, we have
%     \begin{equation}
%         \|\bar F\circ H_{\alpha,\delta}-F\|_{L^p(K)}\le \frac{\varepsilon}{4}+\frac{\varepsilon}{4}+\frac{\varepsilon}{4}+\frac{\varepsilon}{4}=\varepsilon,
%     \end{equation}
%     which completes the proof.
% \end{proof}

\subsection{Proof of \Cref{prop:analytic}}
\label{sec:analytic_proof}
\begin{proof}
Assume that Condition 2 in \Cref{thm:main} fails. Then, there exists $N$ samples $\{X_i\}_{i=1}^N$ from different orbits under the $G$-action, but for any $m$ and $g\in (\operatorname{Id}+\mathcal{G})^m$, there exist indices $i,j\in [N]$ such that at least one token in $g(X_i)$ is identical to a token in $g(X_j)$. This can be written as
\begin{equation}
    \label{eq:fail_}
    \Pi_{i,j}:=\prod_{l_1,l_2} \|[g(X_i)]_{l_1}-[g(X_j)]_{l_2}\|_2^2=0.
\end{equation}
When $\mathcal{G}$ is parametric in $\theta\in\Theta$, $\Pi_{i,j}$ is also analytic in $\theta$.
As the zero set of a nonzero real analytic function has measure zero~\cite{mityagin2015zero}, this implies that for some $i,j$, $\Pi_{i,j}$ is identically zero, meaning that $\mathcal{G}$ fails to distinguish tokens in $X_i$ and $X_j$. This argument reduces Condition 2 in \Cref{thm:main} to the case $N=2$.

Moreover, if there exists a uniform $m$ such that the token distinguishability condition for any two tokens holds, the above argument essentially shows that this $m$ can also be used to distinguish any $N$ tokens.
\end{proof}

\section{Proofs for the applications of main results}
\subsection{Proof of \Cref{thm:kernel} and \Cref{prop:sparse}}
We first prove the following lemma.


% \begin{lemma}[For softmax-type attention \cjp{Need to rewrite for kernel-based attention}]
%     \label{lem:1}
% Suppose $a_1, a_2,\cdots a_r\in\mathbb R^d\setminus \{0\}$ and $b_1,\cdots, b_s\in\mathbb R^d\setminus \{0\}$ are two sequences of different tokens. $m_1,\cdots, m_r$, $n_1,\cdots$ are positive integers. Suppose that for some $r^\prime\in [r]$ and $s^\prime\in [s]$ such that for any $W_Q, W_K$ and $W_V$,

% \begin{equation}
%     \label{eq:lem_equal_atten}
%     \sum_{j=1}^r\left(\frac{m_je^{\left\langle W_Q a_{r^\prime}, W_K a_j\right\rangle}}{\sum_{l=1}^r m_le^{\left\langle W_Q a_{r^\prime}, W_K a_{l}\right\rangle}} W_V a_j\right)= \sum_{j=1}^s \left(\frac{n_je^{\left\langle W_Q b_{s^\prime}, W_K b_j\right\rangle}}{\sum_{l=1}^s n_l e^{\left\langle W_Q b_{s^\prime}, W_K b_{l}\right\rangle}} W_V b_j\right)
% \end{equation}

% Then, we have:
% \begin{enumerate}
%     \item $r=s$, and there exists a permutation $\sigma\in {S}_n$ with $\sigma(r^\prime)=s^\prime$ such that $x_{i}=y_{\sigma(i)}$ for all $i=1,\cdots; n$.
%     \item there exists $q\in \mathbb Q$ such that $m_i=qn_{\sigma(i)}$ for all $i=1,\cdots, r$.
% \end{enumerate}

% \end{lemma}



\begin{lemma}
    \label{lem:equal_atten}
    Let $r$ be a kernel function satisfying the conditions in \Cref{thm:kernel}.
    Suppose $\{a_1, \dots, a_r\}\subset \mathbb{R}^d\setminus \{0\}$ and $\{b_1,\dots, b_s\}\subset \mathbb{R}^d\setminus \{0\}$ are two sequences of distinct tokens. Suppose that for some indices $r^\prime\in [r]$ and $s^\prime\in [s]$, the following equality holds for all choices of $W_Q, W_K, W_V$:
    \begin{equation}
        \label{eq:lem_equal_atten}
        a_{r^\prime}+\sum_{j=1}^r\left(\frac{ k(W_Q a_{r^\prime}, W_K a_j)}{\sum_{l=1}^r k(W_Q a_{r^\prime}, W_K a_l)} W_V a_j\right) =
        b_{s^\prime}+\sum_{j=1}^s\left(\frac{ k(W_Q b_{s^\prime}, W_K b_j)}{\sum_{l=1}^s k(W_Q b_{s^\prime}, W_K b_l)} W_V b_j\right).
    \end{equation}
    Then, $r=s$, and there exists a permutation $\sigma\in{S}_n$ with $\sigma(r^\prime)=s^\prime$ such that $a_i=b_{\sigma(i)}$ for all $i$.
\end{lemma}

\begin{proof}
    Taking $W_V$ to be zero directly gives $a_{r^\prime}=b_{s^\prime}$. In the following, we set $W_Q=I$, and for notational simplicity define
    \[
    \tilde k(x):=k(a_{r^\prime},x),\quad \forall\,x\in\mathbb{R}^d.
    \]
     The proof utilizes the following lemma:
\begin{lemma}[Auxiliary lemma]
    \label{lem:lemma_aux}
    Let $\{x_1,\cdots, x_p\}:=\{a_1,\cdots, a_r\}\cup \{b_1,\cdots, b_s\}$. Then, there exist $W_K\in\mathbb{R}^{d\times d}$ and a permutation $\sigma\in S_p$ such that
    \begin{equation}
        \lim_{t\to\infty} \frac{\tilde k(tW_K x_{\sigma(j)})}{\tilde k(tW_K x_{\sigma(i)})}=\infty, \text{ for all } i<j.
    \end{equation}
\end{lemma}

Since $\{a_j\}_{j=1}^r$ are $r$ distinct tokens, by the auxiliary lemma, we may choose a $W_K$ (and reindex the sequences accordingly) so that after replacing $W_K$ by $tW_K$ the kernel values satisfy, for large $t>0$,
\[
\tilde k(tW_Ka_1)\ll \tilde k(tW_Ka_2)\ll\cdots\ll \tilde k(tW_Ka_r),
\]
and similarly
\[
\tilde k(tW_Kb_1)\ll \tilde k(tW_Kb_2)\ll\cdots\ll \tilde k(tW_Kb_s).
\]
Moreover, the lemma tells us that for any $a_j$ and $b_l$, either $a_j=b_l$ or one of $\tilde k(tW_ka_j)$ and $\tilde k(tW_kb_l)$ is dominated by the other in the limit $t\to\infty$.


After subtracting $a_{r^\prime}$ from both sides in~\eqref{eq:lem_equal_atten}, we have
\begin{equation}
    \label{eq:mean_eq}
    \frac{\sum_{j=1}^r \tilde k(tW_Ka_j)a_j}{\sum_{j=1}^r \tilde k(tW_Ka_j)}
    = \frac{\sum_{j=1}^s \tilde k(tW_Kb_j)b_j}{\sum_{j=1}^s \tilde k(tW_Kb_j)}.
\end{equation}
By a transformation, it gives
\begin{equation}
    \label{eq:mean_eq2}
\sum_{j=1}^r\sum_{l=1}^s \tilde k(tW_Ka_j)\tilde k(tW_Kb_l)(a_j-b_l)=0.
\end{equation}

Let $t\to \infty$, considering the dominate term $\tilde k(tW_Ka_j)\tilde k(tW_Kb_l)(a_r-b_s)$ gives $a_r=b_s$. Then, for all $q=0, 1,\cdots, \min\{r,s\}$, we then prove by induction that: $a_{r-k}=b_{s-k}$.

Suppose we already have $a_{r-i}=b_{s-i}$, for $i=1,\cdots, q-1$. It then follows that

\begin{equation}
        \label{eq:lem_sum_exp_zero}
        \begin{aligned}
                        &\sum_{j=r-q+1}^r\sum_{l=s-q+1}^s \tilde k(tW_Ka_j)\tilde k(tW_Kb_l)(a_j-b_l)\\
                        &=\sum_{j=r-q+1}^r\sum_{l=r-q+1}^r \tilde k(tW_Ka_j)\tilde k(tW_Ka_l)(a_j-a_l)=0, \quad \text{for all }t\in\mathbb R.
        \end{aligned}
\end{equation}.

Combining with~\eqref{eq:mean_eq2}, we have
\begin{equation}
    (\sum_{j=1}^{r-q}\sum_{l=1}^{r-q}+\sum_{j=1}^{r-q}\sum_{l=s-q+1}^s+\sum_{j=r-q+1}^{r}\sum_{l=1}^{s-q}) \left(\tilde k(tW_Ka_j)\tilde k(tW_Kb_l)(a_j-b_l)\right)=0,
\end{equation}
where the leading term is
\begin{equation}
    \tilde k(tW_Ka_{r-q})\tilde k(tW_Kb_{s})(a_{r-q}-b_{s})+\tilde k(tW_Ka_{r})\tilde k(tW_Kb_{s-q})(a_{r}-b_{s-q}).
\end{equation}
Since $a_r=b_s$, $a_{r-q}\neq a_r$ and $b_{r-q}\neq b_r$, let $t\to\infty$ gives that
$\tilde k(tW_Ka_{r-q})$ and $\tilde k(tW_Kb_{s-q})$ are not dominated by each other. By the auxiliary lemma, this indicates that $a_{r-q}=b_{s-q}$, which completes the induction.

Then, we have shown that $a_{r-k}=b_{s-k}$ for all $k=0, 1,\cdots, \min\{r,s\}$. The only remaining thing is to show that $r=s$. Suppose $r<s$, then we have
\begin{equation}
    \sum_{j=1}^r\sum_{l=1}^{s-r} \tilde k(tW_Ka_j)\tilde k(tW_Kb_l)(a_j-b_l)\equiv 0,
\end{equation}
where the unique leading term is $\tilde k(tW_Ka_r)\tilde k(tW_Kb_{s-r})(a_r-b_{s-r})$. Since $a_r=b_r\neq b_{s-r}$, that gives a contradiction.
This completes the proof.
\end{proof}

\begin{proof}[Proof of the auxiliary lemma]
    For each pair $(i,j)$ with \(1 \le i < j \le p\), we define

    \begin{equation}
        \mathcal{M}_{i,j} := \mathbb R^{d\times d}\setminus \left\{ W_K \,\middle|\, \lim_{t \to \infty}\frac{\tilde{k}(tW_Kx_i)}{\tilde{k}(tW_Kx_j)} = 0 \text{ or } \infty\right\},
    \end{equation}
    and we define $\mathcal M$ as the union of all such sets:
    \begin{equation}
        \mathcal{M} = \bigcup_{1 \le i < j \le p} \mathcal{M}_{i,j}.
    \end{equation}
According to the condition in \Cref{thm:kernel}, each $\mathcal M_{i,j}$ is a measure-zero set in $\mathbb R^{d\times d}$. Therefore, $\mathcal M$ is also measure-zero.
Choose any $W_K \in \mathbb R^{d\times d}\setminus \mathcal M$. Then, for any $i,j$, we have
\begin{equation}
    \lim_{t\to\infty} \frac{\tilde k(tW_K x_i)}{\tilde k(tW_K x_j)} = 0 \text{ or } \infty,
\end{equation}
i.e. either $\tilde k(tW_Kx_i)\ll \tilde k(tW_Kx_j)$ or $\tilde k(tW_Kx_i)\gg \tilde k(tW_Kx_j)$ for large $t$. This indicates that there exists a permutation $\sigma\in S_p$ such that
\begin{equation}
    \lim_{t\to\infty} \frac{\tilde k(tW_K x_{\sigma_1(j)})}{\tilde k(tW_K x_{\sigma_1(i)})}=\infty, \text{ for all } i<j,
\end{equation}
which completes the proof.
\end{proof}

\begin{proof}[\textbf{Proof of \Cref{thm:kernel}}]
    \Cref{lem:equal_atten} shows the token-distinguishability condition over $X$ with non-zero tokens.
    This indicates the interpolation property over the region when $X$ has non-zero tokens.
    Since the set
    \begin{equation}
        \{X\in\mathbb R^{d\times n}\mid [X]_i=0 \text{ for some } i\in [n]\}
    \end{equation}
    is a measure-zero set, the $S_n$-UAP holds for $\mathcal T_{\mathcal G,\mathcal H}$ by the same argument as \Cref{thm:main}.
\end{proof}


\subsection{Verification of conditions in \Cref{thm:kernel} on practical attention mechanisms}
\label{sec:verify}
We verify the conditions in \Cref{thm:kernel} for the following kernel functions:
\begin{itemize}
    \item $k(x,y)=\exp(x^\top y)$, used in the original transformer~\cite{vaswani2017attention}.

    \textbf{Verification}: 
    For any given $x\neq 0$ and distinct  $y_1,y_2\in\mathbb R^d\setminus \{0\}$, the set
    \begin{equation}
        \mathcal P:=\{W_K\mid x^\top W_Ky_1=x^\top W_Ky_2\}
    \end{equation}
    is a hyperplane in $\mathbb R^{d\times d}$, which has zero measure. Notice that for any $W_K$ in $\mathbb R^{d\times n}\setminus \mathcal P$,
    \begin{equation}
        \lim_{t\to\infty} \frac{\exp(tx^\top W_Ky_1)}{\exp(tx^\top W_Ky_2)}=\lim_{t\to\infty} \exp(t(x^\top W_K(y_1-y_2)))=\infty \text{ or } 0.
    \end{equation}
    This indicates that the condition in \Cref{thm:kernel} holds.
    \item $k(x,y)=\exp(-\gamma \|x-y\|_2^2)$, the RBF kernel, explored in~\cite{tsai2019transformer}.

    \textbf{Verification:} Notice that
    \begin{equation}
        \frac{\exp(-\gamma\|x-tW_Ky_1\|_2^2)}{\exp(-\gamma\|x-tW_Ky_2\|_2^2)}=\exp(-\gamma(\|tW_Ky_1\|_2^2-\|tW_Ky_2\|_2^2)+2tx^\top W_K (y_2-y_1)).
    \end{equation}

    Therefore, any $W_K$ such that $\|W_Ky_1\|_2\neq \|W_Ky_2\|_2$ satisfies the condition in \Cref{thm:kernel}. Since for distinct $y_1$ and $y_2$,
\begin{equation}
    \|W_Ky_1\|_2^2-\|W_Ky_2\|_2^2= 0,
\end{equation}
is a non-zero quadratic equation on $W_K$, whose solution set has zero measure. Therefore, the condition in \Cref{thm:kernel} holds.

    \item $k(x,y)=\phi(x)^\top\phi(y)$, where
    \begin{equation}
        \phi(x)^{\top}=\exp \left(-\frac{1}{2}\|x\|^2\right)\left(\exp ({\omega}_1^{\top} x), \ldots, \exp ({\omega}_m^{\top} x)\right)\in\mathbb{R}^m,
    \end{equation}
    with $\omega_1,\dots, \omega_m\in\mathbb{R}^d$ drawn i.i.d. from a Gaussian distribution. This kernel is used in Performer~\cite{choromanski2020rethinking}, and \Cref{thm:kernel} holds almost surely in this case.

    \textbf{Verification:}
    We have
    \begin{equation}
        \label{eq:kernel_phi}
        \begin{aligned}
            \frac{k(x,tW_Ky_1)}{k(x,tW_Ky_2)}&=\frac{\phi(x)^{\top}\phi(tW_Ky_1)}{\phi(x)^{\top}\phi(tW_Ky_2)}\\
            &=\exp\left(\frac{t^2}{2}(\|W_Ky_2\|_2^2-\|W_Ky_1\|_2^2)\right)\frac{\sum_{i=1}^m \exp(\omega_i^\top x) \exp( t \omega_i^\top W_ky_1)}{\sum_{i=1}^m \exp(\omega_i^\top x)\exp(t \omega_i^\top W_ky_2)}
        \end{aligned}
    \end{equation}
    We claim that if $w_i$ are pair-wise linear independent, i.e. there does not exist $i\neq j$ such that $w_i=\alpha w_j$ for some $\alpha\in\mathbb R$, the condition in \Cref{thm:kernel} holds. This almost surely holds when $w_i$ are drawn i.i.d. from a Gaussian distribution.

    If when $t\to\infty$, the ratio in~\eqref{eq:kernel_phi} do not goes to infinity or zero, it must hold that
    \begin{equation}
        \label{eq:kernel_phi2}
        \|W_Ky_1\|_2=\|W_Ky_2\|_2, \text{ and }         \max_{i}\{\omega_i^\top W_Ky_1\}=\max_{i}\{\omega_i^\top W_Ky_2\}.
    \end{equation}
    When $w_i$ are pair-wise linear independent, we have that $y_1w_i^\top\neq y_2w_j^\top$ for all $i\neq j\in [m]$. Also, since $y_1\neq y_2$ and $w_i$ are non-zero(by the pair-wise independent condition), we have $y_1w_i^\top\neq y_2w_i^\top$ for all $i\in [m]$. Therefore, we have that all the sets
    \begin{equation}
        \{W_K\mid \omega_i^\top W_Ky_1=\omega_j^\top W_Ky_1\}=\{W_K\mid \langle W_K, y_1 w_i^\top-y_2w_j^\top\rangle_F=0\},
    \end{equation}
    where $\langle\cdot, \cdot\rangle_F$ denotes the Frobenius inner product, are hyperplanes in $\mathbb R^{d\times d}$, which has zero measure. That is, equation~~\eqref{eq:kernel_phi2} only holds for a measure-zero set of $W_K$, which completes the verification.
    \item $k(x,y)=\exp(w^\top x)+\exp(w^\top y)$, where $w\in\mathbb R^d$.
 
    \textbf{Verification:}
    We have
    \begin{equation}
        \frac{k(x,tW_Ky_1)}{k(x,tW_Ky_2)}=\frac{\exp(w^\top (x+ tW_Ky_1))}{\exp(w^\top (x+tW_Ky_2))}=\exp(tw^\top W_K(y_1-y_2)).
    \end{equation}
    When $y_1\neq y_2$, for almost all $W_K$, we have $w^\top W_K(y_1-y_2)\neq 0$. Therefore, the condition in \Cref{thm:kernel} holds.
    \item $k(x,y)=p(x-y)\tilde k(x,y)$, with $p$ being any positive polynomial function and $\tilde k$ being any kernel satisfies the condition in \Cref{thm:kernel}. 
    
    \textbf{Verification:} Just neet to notice that for almost all $W_k$, it holds that
    \begin{equation}
        \lim_{t\to\infty} \frac{p(x-tW_Ky_1)}{p(x-ptW_Ky_2)}
    \end{equation}
    is a constant indicating that the condition in \Cref{thm:kernel} still holds.
\end{itemize}

\subsection{Proof of \Cref{prop:sparse}}
\begin{proof}
    We only need to prove the token distinguishability condition for two samples:
    \begin{itemize}
        \item For any $X$ and $Y$ that are in general positions and from different orbits of $G$(defined in~\eqref{eq:sparse_group}), there exists
    \begin{equation}
     g\in \mathcal G_m^\Phi:=\{(\operatorname{Id}+g_m)\circ \cdots \circ (\operatorname{Id}+g_1)\mid g_i\in\mathcal G_{\mathcal N_i}, \text{ for } i \in [m]\}
    \end{equation}
     such that the tokens of $g(X)$ and $g(Y)$ are all distinct.
    \end{itemize}
    We prove this claim by contradiction. Assume that, there exist $X$ and $Y$ that are in general positions and from different orbits of $G$, but for any $g\in \mathcal G_m^\Phi$, there exist indices $i,j\in [n]$ such that at least one token in $g(X)$ is identical to a token in $g(Y)$.
    Then, according to the analyticity, there exist indices $i_1$ and $i_2$ such that, $[g(X)]_{i_1}=[g(Y)]_{i_2}$ always hold. For a given $p_1\in [m]$, we first consider $g\in\mathrm{Id} +\mathcal G_{\mathcal N_{p_1}}$. Then, $[g(X)]_{i_1}=[g(Y)]_{i_2}$ gives:
    \begin{equation}
        \label{eq:prop2_equal_atten}
[X]_{i_1}+\frac{\sum_{j\in\mathcal N_{p_1}(i_1)} k([W_QX]_{i_1},[W_KX]_j)[W_VX]_j}{\sum_{j\in\mathcal N_{p_1}(i_1)} k([W_QX]_{i_1},[W_KX]_j)}= [Y]_{i_2}+\frac{\sum_{l\in\mathcal N_{p_1}(i_2)} k([W_QY]_{i_2},[W_KY]_l)[W_VY]_l}{\sum_{j\in\mathcal N_{p_1}(i_2)} k([W_QY]_{i_2},[W_KY]_l)},
    \end{equation}
    for any $W_Q, W_K, W_V\in\mathbb{R}^{d\times d}$. According to \Cref{lem:equal_atten}, we can deduce that $|\mathcal N_{p_1}(i)|=|\mathcal N_{p_1}(j)|$, $[X]_{i_1}=[Y]_{i_2}$, and 
    % the set of tokens with indices in $\mathcal N_{p_1}(i_1)$ of $X$ and those in $\mathcal N_{p_1}(i_2)$ of $Y$ are the same. That shows that for any $q_1\in \mathcal N(i_1)$, there exists $q_2\in \mathcal N(i_2)$ such that $[X]_{q_1}=[Y]_{q_2}$, i.e. 
    % the set of tokens of $X$ that $[X]_{i_1}$ can directly attend to is the same as those of $Y$ that $[Y]_{i_2}$ can directly attend to.
    \begin{equation}
        \{[X]_q\mid q\in \mathcal N_{p_1}(i_1)\}=\{[Y]_q\mid q\in \mathcal N_{p_1}(i_2)\}.
    \end{equation}
    
    Therefore, choose any $q_1\in \mathcal N_{p_1}(i_1)$, we can find $q_2\in \mathcal N_{p_1}(i_2)$ such that $[X]_{q_1}=[Y]_{q_2}$. 

    Now, we claim that for any $p_2<p_1$
    and $g\in \mathrm{Id}+\mathcal G_{\mathcal N_{p_2}}$(a layer before the $p_1$-th layer) , $[g(X)]_{q_1}=[g(Y)]_{q_2}$.
    Otherwise, suppose there exists $g_1\in \mathrm{Id}+\mathcal G_{p_2}$ with $[g_1(X)]_{q_1}\neq [g_1(Y)]_{q_2}$.
    By scaling the $W_V$ matrix to be small enough, we can assume $g_1$ satisfies that
    \begin{equation}
        \label{eq:prop2_distinct}
        \|g_1(X)-X\|_2<\frac{1}{2}\min_{i\neq j}\{\|[X]_i-[X]_j\|_2\}.
    \end{equation}
Then, we have that $[g_1(X)]_{q_1}\neq [g_1(Y)]_{q_2}$, and by equation~\eqref{eq:prop2_distinct}, for any $q\neq q_2$ in $\mathcal N(i_2)$, we have
\begin{equation}
    \begin{aligned}
        \|[g_1(X)]_{q_1}-[g_1(Y)]_q\|_2&=  \|([g_1(X)]_{q_1}-[X]_{q_1})+([X]_{pq1}-[Y]_q)+([Y]_q-[g_1(Y)]_q)\|_2\\
        &\ge \|[X]_{g_1}-[Y]_q\|_2-\|[q_1(X)]_{q_1}-[X]_{q_1}\|_2-\|[Y]_q-[g_1(Y)]_q\|_2>0.
    \end{aligned}
\end{equation}
Therefore, $[g_1(X)]_{q_1}$ does not appear in the tokens of $g_1(Y)$, i.e. the sets
\begin{equation}
    \{[g_1(X)]_j\mid j\in\mathcal N(i_1)\}\quad\text{ and } \quad \{[g_1(Y)]_l\mid l\in\mathcal N(i_2)\}
\end{equation}
must be different. By applying \Cref{lem:equal_atten} to $[g_1(X)]_{q_1}$ and $[g_1(Y)]_{q_2}$, we know that there exists $g_2\in \mathrm{Id}+\mathcal G_{\mathcal N_{p_1}}$ such that $[g_2(g_1(X))]_{q_1}\neq [g_2(g_1(Y))]_{q_2}$. Since $g_2\circ g_1\in (\mathrm{Id}+\mathcal G_{\mathcal N_{p_2}})\circ (\mathrm{Id}+\mathcal G_{\mathcal N_{p_1}})\subset \mathcal G_m^\Phi$ which can distinguish $[X]_{i_1}$ and $[Y]_{i_2}$, contradicting to our assumption.


Hence, we have shown that for $p_2< p_1\le n$ and any $g \in \mathrm{Id}+\mathcal G_{\mathcal N_{p_2}}$, $[g(X)]_{q_1}=[g(Y)]_{q_2}$. Then, we can apply \Cref{lem:equal_atten} to $[X]_{q_1}$ and $[Y]_{q_2}$, and deduce that the set of tokens of $X$ with indices in $\mathcal N_{p_s}(q_1)$ are the same as those of $Y$ with indices in $\mathcal N_{p_3}(q_2)$. That is, the tokens where $[X]_{i_1}$ can attend to within two hops are the same as those where $[Y]_{i_2}$ can attend to within two hops.

The above process can be repeated. Since we assume that $\Phi$ is connected within $m$ layers, we know that any indices in $[n]$ can be reached starting from $i_1$ and $i_2$ within $m$ hops. Finally, the above discussion can cover all indices in $[n]$. Since $X$ and $Y$ are in general positions, the correspondence between their tokens is unique.
Finally, this results in a permutation $\sigma\in S_n$, such that:
\begin{equation}
    [X]_i=[Y]_{\sigma(i)}.
\end{equation}
On the other hand, apply~\Cref{lem:equal_atten} again to eqch $\mathcal N_p$, $[X]_i$ and $[Y]_{\sigma(i)}$, we can deduce that
\begin{equation}
     j\in\mathcal N_p(i) \Leftrightarrow \sigma(j)\in\mathcal N_p(\sigma(i)), \text{ for all } p\in [m].
\end{equation}
By the definition of $\operatorname{Aut}(\mathcal N)$, this means that $\sigma$ belongs to each $\mathcal N_{ p}$, indicating that $\sigma \in G$.
However, this contradicts to our assumption that $X$ and $Y$ are from different orbits of $\operatorname{Aut}(\Gamma)$, which completes the proof.
\end{proof}

\subsection{Verification of the UAP for other transformer variants}
\label{sec:other_transformer}
In this section, we verify the condition in \Cref{thm:main} for the kernelized attention of SkyFormer~\cite{chen2021skyformer} and the attention mechanism of the Linformer~\cite{wang2020linformer}.

\subsubsection{UAP for LinFormer}
Linformer~\cite{wang2020linformer}, where the attention layer is defined as
    \begin{equation}
        % \label{eq:linformer}
        \operatorname{Atten}(X)=X + W_V X {F}\operatorname{softmax}((W_K X{E})^\top W_Q X)
\end{equation}
where $E, F\in \mathbb R^{n\times k}$ with $1\le k\ll n$ are two trainable projection matrices.

For LinFormer, we have the following lemma:
\begin{lemma}
    Let $X,Y\in \mathbb R^{d\times n}$ be two points that are in general positions. If for some $i_1,i_2\in [n]$, the following equality holds for all $W_Q, W_K, W_V\in\mathbb R^{d\times d}$ and $E, F\in\mathbb R^{n\times k}$:
    \begin{equation}
        \label{eq:lem_equal_atten_linformer}
        [X]_{i_1}+[W_V X {F}\operatorname{softmax}((W_K X{E})^\top W_Q X)]_{i_1} = [Y]_{i_2}+[W_V Y {F}\operatorname{softmax}((W_K Y{E})^\top W_Q Y)]_{i_2},
    \end{equation}
     then we have $i_1=i_2$ and $X=Y$.
\end{lemma}

\begin{proof}
    Take $W_V=0$ gives $[X]_{i_1}=[Y]_{i_2}$. Then, for any given $R=\{r_1,\cdots, r_k\}\subset [n]$, we take
    \begin{equation}
        E=F=[e_{r_1},\cdots, e_{r_k}],
    \end{equation}
    where $e_{r_i}$ is the $r_i$-th column of the identity matrix. Equation~\eqref{eq:lem_equal_atten_linformer} then gives that
    \begin{equation}
        \sum_{j\in R}\left(\frac{ \exp(\langle W_Q [X]_{i_1}, W_K [X]_j\rangle)}{\sum_{l\in R} \exp(\langle W_Q [X]_{i_1}, W_K [X]_l\rangle)} W_V [X]_j \right) =          \sum_{j\in R}\left(\frac{ \exp(\langle W_Q [Y]_{i_2}, W_K [Y]_j\rangle)}{\sum_{l\in R} \exp(\langle W_Q [Y]_{i_2}, W_K [Y]_l\rangle)} W_V [Y]_j \right)
    \end{equation}
    which reduces to the discussion in \Cref{lem:equal_atten}. Therefore, we have that the set $\{[X]_{i}\mid i\in R\}$ is the same as $\{[Y]_{i}\mid i\in R\}$. Since $X$ and $Y$ are in general positions, and $R$ is arbitrary, this indicates that $X=Y$.
\end{proof}

According to this lemma and the fact that~\eqref{eq:linformer} is analytic to all the parameters, we conclude by \Cref{prop:analytic} that the UAP holds for LinFormer without symmetric restrictions. Furthermore, the same result can be generalized to the case where the softmax function in~\eqref{eq:linformer} is replaced by a kernel-based form with a kernel satisfying the condition in \Cref{thm:kernel}.

\subsubsection{UAP for SkyFormer}
 The kernelized attention used in SkyFormer~\cite{chen2021skyformer}, where the attention mechanism is given by:
    \begin{equation}
        \label{eq:skyformer}
        [\operatorname{Atten}(X)]_i=[X]_i + \sum_{j=1}^n \exp{\left(-\frac{1}{2}\|[W_QX]_i-[W_kX]_j\|^2\right)}W_VX_j.
    \end{equation}

The proof follows from the verification for the RBF kernel \Cref{sec:verify}, with the same argument to prove the token-distinguishability condition.

\subsubsection{UAP for architecture proposed in~\eqref{eq:new_atten}}
\label{sec:proof_new_atten}
More precisely, if we define $\tilde{\mathcal G}_{\mathcal N}$ as the family of token-mixing maps associated with the sparsity pattern $\mathcal N$, and the transformer family $\tilde{\mathcal T}_{\mathcal H}^\Phi$ associated with a sequence of sparse mode, just as $\mathcal T_{\mathcal H}^\Phi$ defined in~\Cref{sec:sparse_atten}. Then, under the same assumption on $\Phi$ as in~\Cref{prop:sparse}, we have:
\begin{corollary}
    \label{thm:new_atten}
    % Suppose that $\Phi$ is connected within $m$ layers.
\end{corollary}
$\tilde {\mathcal T}_{\mathcal H}^\Phi$ possesses the $G$-UAP with $G$ defined in~\eqref{eq:sparse_group}.

Assume the condition in \Cref{prop:analytic} fails. Then, there exists $X$ and $Y$ that are in general positions, and $i_1, i_2\in [n]$ such that for all $W\in\mathbb R^{d\times d}$, $a\in \mathbb R$ and $b\in \mathbb R^d$, we have
\begin{equation}
    [X]_{i_1}+\sum_{j\in\mathcal N(i_1)} a \alpha(W[X]_j-b)=[Y]_{i_2}+\sum_{j\in\mathcal N(i_2)} a \alpha(W[Y]_j-b).
\end{equation}
It then follows that $[X]_{i_1}=[Y]_{i_2}$. Moreover, if the sets 
\begin{equation}
    \{[X]_j\mid j\in \mathcal N(i_1)\}
\end{equation}
and 
\begin{equation}
    \{[Y]_j\mid j\in \mathcal N(i_2)\}
\end{equation}
are not the same, we can then derive an identity
\begin{equation}
    \label{eq:new_atten_equal}
    \sum_{l=1}^L c_l\alpha(Wx_l-b)=0,
\end{equation}
where $x_l$ are the unique tokens appears in the two sets, $c_i$ equals to $1$ if $x_l$ appears in the set of $X$, and $-1$ otherwise. Notice that the identity holds for all $W\in\mathbb R^{d\times d}$ and $b\in\mathbb R$, this gives a contradiction according to the proof of Theorem 1 in~\cite{cheng2025interpolation}.
Specifically, since $\alpha$ is of polynomial growth, taking the Fourier transform(in distributional sense) on both side of~\eqref{eq:new_atten_equal} with respect to $b$ gives:
\begin{equation}
    (\sum_{l=1}^L c_l e^{i (Wx_l)^\top \xi})\hat \alpha(\xi)=0, \quad \text{ for all } \xi\in\mathbb R^d,
\end{equation}
where $\hat \alpha$ is the Fourier transform of $\alpha$. 
Since $\alpha$ is not a polynomial, we have $\operatorname{supp} \hat \alpha$ contains a non-zero value.
This will lead to a contradiction. See Section 3.2 in~\cite{cheng2025interpolation} for more details.

\subsubsection{Details for architecture proposed for $D_n$/$C_n$ equivariant map}
\textbf{For architecture with $D_n$ symmetry}
By choosing $\Phi=(\mathcal N,\mathcal N,\cdots)$ as an invariant sequence of sparse mode, by~\Cref{prop:sparse}, we have that the transformer with the first design satisfies the $\operatorname{Aut}(\mathcal N)$-UAP. Therefore, we only need to prove thet $\operatorname{Aut}(\mathcal N)=D_n$. 

First, it is easy to see that for any $g\in D_n$ , we have $j\in \mathcal N(i)$ indicates that $g(j)\in \mathcal N(g(i))$. This implies that $N_n\in \operatorname{Aut}(\mathcal N)$. On the other hand, for any $g\in \operatorname{Aut}(\mathcal N)$, since $D_n$ is transitive, there exists $h\in D_n$ such that $h(g(1))=1$. Now, we prove that $\gamma:=h\circ g$ is either the identity or the reflection $\sigma=(1, n)(2, n-1)\cdots$.

Since $\gamma\in \operatorname{Aut}(\mathcal N)$, we know that $\gamma(\mathcal N(i))=\mathcal N(\gamma(i))$. Since $\gamma(1)=1$, we have $\gamma(\mathcal N(1))=\mathcal N(\gamma(1))=\mathcal N(1)$ is invariant. This indicates that for any $i\in\mathcal N(1)$, $\gamma(i)$ is also in $\mathcal N(1)$. Now, we consider the value of $\mathcal N(2)$. We have that 
\begin{equation}
    2w=|\mathcal N(1)\cap \mathcal N(2)|=|\gamma(\mathcal N(1))\cap \gamma(\mathcal N(2))|=|\mathcal N(1)\cap \mathcal N(\gamma(2))|
\end{equation}
Since $2w+1\le n-2$, we have that in $\mathcal N(1)$, there are only two indices $j=n,2$ such that 
\begin{equation}
    |\mathcal N(1)\cap \mathcal N(j)|=2w.
\end{equation}
Therefore, it follows that $\gamma(2)=n$ or $2$. If $\gamma(2)=2$, we can then repeat the discussion to deduce that $\gamma(i)=i$ for $i=2,3\cdots, n$ sequentially, indicating that $\gamma$ is the identity. 
If $\gamma(2)=n$. Then, we can repeat the discussion to deduce that $\gamma(i)=n+2-i \operatorname{mod} n$ for $i=2,3\cdots, n$. This indicates that $\gamma$ is a reflection $(2, n)(3,n-1)\cdots$, which is in $D_n$.

Therefore, we have shown that $\operatorname{Aut}(\mathcal N)=D_n$. This is actually a classical result on the automorphism group of the circulant  graph~\cite{biggs1993algebraic}.


Moreover, if we destroy the symmetry to reflection by defining 
\begin{equation}
    \mathcal N(i):=\{i, i+1,i+2,\cdots, i+w \operatorname{mod} n\}\text{ for } i=1,2,\cdots, n,
\end{equation}
with $w\le\lfloor \frac{n-1}{2}\rfloor-1$, we get a transformer that is $C_n$-equivariant and satisfies the $C_n$-UAP. For the proof, we only need to check that $\operatorname{Aut}(\mathcal N)=C_n $, which can be done following the same approach as $D_n$.

\textbf{For architecture with $C_n$ symmetry}
For token-mixing layer defined by the convolution in~\eqref{eq:conv}, we first notice that it satisfies the $C_n$-equivariance. In fact, this follows from the fact that the convolutional operation is equivariant to translation.

Therefore, to prove the $C_n$-UAP, we only need to check the token distinguishability condition under $C_n$ action. Specifically, suppose the condition in~\Cref{prop:analytic} fails. Then, there exists $X$ and $Y$ that are in general positions, and $i_1,i_2\in [n]$ such that the composition of token mixing layers cannot distinguish the $i_1$-th token of $X$ and the $i_2$-th token of $Y$. 

Considering using single layers, we have that for all $\psi\in\mathbb R^{l+1}$, it holds

\begin{equation}
    [\psi * X]_{i_1} = [\psi * X]_{i_2},
\end{equation}
i.e. 
\begin{equation}
    \sum_{j=0}^{l} \psi_{j}\, [X]_{(i_1 + j) \bmod n}=\sum_{j=0}^{l} \psi_{j}\, [Y]_{(i_2 + j) \bmod n},
\end{equation}
which indicates that
\begin{equation}
    \sum_{j=0}^{l} \psi_{j}\, ([X]_{(i_1 + j) \bmod n}-[Y]_{(i_2 + j) \bmod n})=0
\end{equation}
%%%%%%%%%%%%%%%%%%%%%%%%%%%%%%%%%%%%%%%%%%%%%%%%%%%%%%%%%%%%
Since $\psi$ is arbitrary, this indicates that
\begin{equation}
    [X]_{(i_1 + j) \bmod n}=[Y]_{(i_2 + j) \bmod n}, \quad \text{ for } j=0,1,\cdots, l.
\end{equation}
Then, we takethe indicices ${i_1+ l\operatorname{mod} n}$ and ${i_2+ l\operatorname{mod} n}$ of $X, Y$ respectively, and consider using two layers. The process is essentially the same as the proof of \Cref{prop:sparse}. We can finally deduce that 
\begin{equation}
    [X]_{(i_1 + j) \bmod n}=[Y]_{(i_2 + j) \bmod n}, \quad \text{ for } j=0,1,\cdots, n.
\end{equation}
That is, $X$ and $Y$ differs only a cyclic action on tokens, meaning that they are from the same $C_n$ orbit, which is a contradiction, and completes the proof.

\newpage
\section*{NeurIPS Paper Checklist}

%%% BEGIN INSTRUCTIONS %%%
% The checklist is designed to encourage best practices for responsible machine learning research, addressing issues of reproducibility, transparency, research ethics, and societal impact. Do not remove the checklist: {\bf The papers not including the checklist will be desk rejected.} The checklist should follow the references and follow the (optional) supplemental material.  The checklist does NOT count towards the page
% limit. 

% Please read the checklist guidelines carefully for information on how to answer these questions. For each question in the checklist:
% \begin{itemize}
%     \item You should answer \answerYes{}, \answerNo{}, or \answerNA{}.
%     \item \answerNA{} means either that the question is Not Applicable for that particular paper or the relevant information is Not Available.
%     \item Please provide a short (1–2 sentence) justification right after your answer (even for NA). 
%    % \item {\bf The papers not including the checklist will be desk rejected.}
% \end{itemize}

% {\bf The checklist answers are an integral part of your paper submission.} They are visible to the reviewers, area chairs, senior area chairs, and ethics reviewers. You will be asked to also include it (after eventual revisions) with the final version of your paper, and its final version will be published with the paper.

% The reviewers of your paper will be asked to use the checklist as one of the factors in their evaluation. While "\answerYes{}" is generally preferable to "\answerNo{}", it is perfectly acceptable to answer "\answerNo{}" provided a proper justification is given (e.g., "error bars are not reported because it would be too computationally expensive" or "we were unable to find the license for the dataset we used"). In general, answering "\answerNo{}" or "\answerNA{}" is not grounds for rejection. While the questions are phrased in a binary way, we acknowledge that the true answer is often more nuanced, so please just use your best judgment and write a justification to elaborate. All supporting evidence can appear either in the main paper or the supplemental material, provided in appendix. If you answer \answerYes{} to a question, in the justification please point to the section(s) where related material for the question can be found.

% IMPORTANT, please:
% \begin{itemize}
%     \item {\bf Delete this instruction block, but keep the section heading ``NeurIPS Paper Checklist"},
%     \item  {\bf Keep the checklist subsection headings, questions/answers and guidelines below.}
%     \item {\bf Do not modify the questions and only use the provided macros for your answers}.
% \end{itemize} 
 

%%% END INSTRUCTIONS %%%


\begin{enumerate}

\item {\bf Claims}
    \item[] Question: Do the main claims made in the abstract and introduction accurately reflect the paper's contributions and scope?
    \item[] Answer: \answerYes{} % Replace by \answerYes{}, \answerNo{}, or \answerNA{}.
    \item[] Justification: The claims made in the abstract and introduction accurately reflect the paper's contributions and scope. 
    \item[] Guidelines:
    \begin{itemize}
        \item The answer NA means that the abstract and introduction do not include the claims made in the paper.
        \item The abstract and/or introduction should clearly state the claims made, including the contributions made in the paper and important assumptions and limitations. A No or NA answer to this question will not be perceived well by the reviewers. 
        \item The claims made should match theoretical and experimental results, and reflect how much the results can be expected to generalize to other settings. 
        \item It is fine to include aspirational goals as motivation as long as it is clear that these goals are not attained by the paper. 
    \end{itemize}

\item {\bf Limitations}
    \item[] Question: Does the paper discuss the limitations of the work performed by the authors?
    \item[] Answer: \answerYes{} % Replace by \answerYes{}, \answerNo{}, or \answerNA{}.
    \item[] Justification: We discuss the limitations of our work in the last section of the paper.
    \item[] Guidelines:
    \begin{itemize}
        \item The answer NA means that the paper has no limitation while the answer No means that the paper has limitations, but those are not discussed in the paper. 
        \item The authors are encouraged to create a separate "Limitations" section in their paper.
        \item The paper should point out any strong assumptions and how robust the results are to violations of these assumptions (e.g., independence assumptions, noiseless settings, model well-specification, asymptotic approximations only holding locally). The authors should reflect on how these assumptions might be violated in practice and what the implications would be.
        \item The authors should reflect on the scope of the claims made, e.g., if the approach was only tested on a few datasets or with a few runs. In general, empirical results often depend on implicit assumptions, which should be articulated.
        \item The authors should reflect on the factors that influence the performance of the approach. For example, a facial recognition algorithm may perform poorly when image resolution is low or images are taken in low lighting. Or a speech-to-text system might not be used reliably to provide closed captions for online lectures because it fails to handle technical jargon.
        \item The authors should discuss the computational efficiency of the proposed algorithms and how they scale with dataset size.
        \item If applicable, the authors should discuss possible limitations of their approach to address problems of privacy and fairness.
        \item While the authors might fear that complete honesty about limitations might be used by reviewers as grounds for rejection, a worse outcome might be that reviewers discover limitations that aren't acknowledged in the paper. The authors should use their best judgment and recognize that individual actions in favor of transparency play an important role in developing norms that preserve the integrity of the community. Reviewers will be specifically instructed to not penalize honesty concerning limitations.
    \end{itemize}

\item {\bf Theory assumptions and proofs}
    \item[] Question: For each theoretical result, does the paper provide the full set of assumptions and a complete (and correct) proof?
    \item[] Answer: \answerYes{} % Replace by \answerYes{}, \answerNo{}, or \answerNA{}.
    \item[] Justification: We provide exact and complete assumptions and proofs for all theoretical results in the paper and the appendix.
    \item[] Guidelines:
    \begin{itemize}
        \item The answer NA means that the paper does not include theoretical results. 
        \item All the theorems, formulas, and proofs in the paper should be numbered and cross-referenced.
        \item All assumptions should be clearly stated or referenced in the statement of any theorems.
        \item The proofs can either appear in the main paper or the supplemental material, but if they appear in the supplemental material, the authors are encouraged to provide a short proof sketch to provide intuition. 
        \item Inversely, any informal proof provided in the core of the paper should be complemented by formal proofs provided in appendix or supplemental material.
        \item Theorems and Lemmas that the proof relies upon should be properly referenced. 
    \end{itemize}

    \item {\bf Experimental result reproducibility}
    \item[] Question: Does the paper fully disclose all the information needed to reproduce the main experimental results of the paper to the extent that it affects the main claims and/or conclusions of the paper (regardless of whether the code and data are provided or not)?
    \item[] Answer: \answerNA{} % Replace by \answerYes{}, \answerNo{}, or \answerNA{}.
    \item[] Justification: This paper does not include experiments.
    \item[] Guidelines:
    \begin{itemize}
        \item The answer NA means that the paper does not include experiments.
        \item If the paper includes experiments, a No answer to this question will not be perceived well by the reviewers: Making the paper reproducible is important, regardless of whether the code and data are provided or not.
        \item If the contribution is a dataset and/or model, the authors should describe the steps taken to make their results reproducible or verifiable. 
        \item Depending on the contribution, reproducibility can be accomplished in various ways. For example, if the contribution is a novel architecture, describing the architecture fully might suffice, or if the contribution is a specific model and empirical evaluation, it may be necessary to either make it possible for others to replicate the model with the same dataset, or provide access to the model. In general. releasing code and data is often one good way to accomplish this, but reproducibility can also be provided via detailed instructions for how to replicate the results, access to a hosted model (e.g., in the case of a large language model), releasing of a model checkpoint, or other means that are appropriate to the research performed.
        \item While NeurIPS does not require releasing code, the conference does require all submissions to provide some reasonable avenue for reproducibility, which may depend on the nature of the contribution. For example
        \begin{enumerate}
            \item If the contribution is primarily a new algorithm, the paper should make it clear how to reproduce that algorithm.
            \item If the contribution is primarily a new model architecture, the paper should describe the architecture clearly and fully.
            \item If the contribution is a new model (e.g., a large language model), then there should either be a way to access this model for reproducing the results or a way to reproduce the model (e.g., with an open-source dataset or instructions for how to construct the dataset).
            \item We recognize that reproducibility may be tricky in some cases, in which case authors are welcome to describe the particular way they provide for reproducibility. In the case of closed-source models, it may be that access to the model is limited in some way (e.g., to registered users), but it should be possible for other researchers to have some path to reproducing or verifying the results.
        \end{enumerate}
    \end{itemize}


\item {\bf Open access to data and code}
    \item[] Question: Does the paper provide open access to the data and code, with sufficient instructions to faithfully reproduce the main experimental results, as described in supplemental material?
    \item[] Answer: \answerNA{} % Replace by \answerYes{}, \answerNo{}, or \answerNA{}.
    \item[] Justification: This paper does not include experiments.
    \item[] Guidelines:
    \begin{itemize}
        \item The answer NA means that paper does not include experiments requiring code.
        \item Please see the NeurIPS code and data submission guidelines (\url{https://nips.cc/public/guides/CodeSubmissionPolicy}) for more details.
        \item While we encourage the release of code and data, we understand that this might not be possible, so “No” is an acceptable answer. Papers cannot be rejected simply for not including code, unless this is central to the contribution (e.g., for a new open-source benchmark).
        \item The instructions should contain the exact command and environment needed to run to reproduce the results. See the NeurIPS code and data submission guidelines (\url{https://nips.cc/public/guides/CodeSubmissionPolicy}) for more details.
        \item The authors should provide instructions on data access and preparation, including how to access the raw data, preprocessed data, intermediate data, and generated data, etc.
        \item The authors should provide scripts to reproduce all experimental results for the new proposed method and baselines. If only a subset of experiments are reproducible, they should state which ones are omitted from the script and why.
        \item At submission time, to preserve anonymity, the authors should release anonymized versions (if applicable).
        \item Providing as much information as possible in supplemental material (appended to the paper) is recommended, but including URLs to data and code is permitted.
    \end{itemize}


\item {\bf Experimental setting/details}
    \item[] Question: Does the paper specify all the training and test details (e.g., data splits, hyperparameters, how they were chosen, type of optimizer, etc.) necessary to understand the results?
    \item[] Answer: \answerNA{} % Replace by \answerYes{}, \answerNo{}, or \answerNA{}.
    \item[] Justification: This paper does not include experiments.
    \item[] Guidelines:
    \begin{itemize}
        \item The answer NA means that the paper does not include experiments.
        \item The experimental setting should be presented in the core of the paper to a level of detail that is necessary to appreciate the results and make sense of them.
        \item The full details can be provided either with the code, in appendix, or as supplemental material.
    \end{itemize}

\item {\bf Experiment statistical significance}
    \item[] Question: Does the paper report error bars suitably and correctly defined or other appropriate information about the statistical significance of the experiments?
    \item[] Answer: \answerNA{} % Replace by \answerYes{}, \answerNo{}, or \answerNA{}.
    \item[] Justification: This paper does not include experiments.
    \item[] Guidelines:
    \begin{itemize}
        \item The answer NA means that the paper does not include experiments.
        \item The authors should answer "Yes" if the results are accompanied by error bars, confidence intervals, or statistical significance tests, at least for the experiments that support the main claims of the paper.
        \item The factors of variability that the error bars are capturing should be clearly stated (for example, train/test split, initialization, random drawing of some parameter, or overall run with given experimental conditions).
        \item The method for calculating the error bars should be explained (closed form formula, call to a library function, bootstrap, etc.)
        \item The assumptions made should be given (e.g., Normally distributed errors).
        \item It should be clear whether the error bar is the standard deviation or the standard error of the mean.
        \item It is OK to report 1-sigma error bars, but one should state it. The authors should preferably report a 2-sigma error bar than state that they have a 96\% CI, if the hypothesis of Normality of errors is not verified.
        \item For asymmetric distributions, the authors should be careful not to show in tables or figures symmetric error bars that would yield results that are out of range (e.g. negative error rates).
        \item If error bars are reported in tables or plots, The authors should explain in the text how they were calculated and reference the corresponding figures or tables in the text.
    \end{itemize}

\item {\bf Experiments compute resources}
    \item[] Question: For each experiment, does the paper provide sufficient information on the computer resources (type of compute workers, memory, time of execution) needed to reproduce the experiments?
    \item[] Answer: \answerNA{} % Replace by \answerYes{}, \answerNo{}, or \answerNA{}.
    \item[] Justification: This paper does not include experiments.
    \item[] Guidelines:
    \begin{itemize}
        \item The answer NA means that the paper does not include experiments.
        \item The paper should indicate the type of compute workers CPU or GPU, internal cluster, or cloud provider, including relevant memory and storage.
        \item The paper should provide the amount of compute required for each of the individual experimental runs as well as estimate the total compute. 
        \item The paper should disclose whether the full research project required more compute than the experiments reported in the paper (e.g., preliminary or failed experiments that didn't make it into the paper). 
    \end{itemize}
    
\item {\bf Code of ethics}
    \item[] Question: Does the research conducted in the paper conform, in every respect, with the NeurIPS Code of Ethics \url{https://neurips.cc/public/EthicsGuidelines}?
    \item[] Answer: \answerYes{} % Replace by \answerYes{}, \answerNo{}, or \answerNA{}.
    \item[] Justification: The research conducted in this paper conforms with the NeurIPS Code of Ethics.
    \item[] Guidelines:
    \begin{itemize}
        \item The answer NA means that the authors have not reviewed the NeurIPS Code of Ethics.
        \item If the authors answer No, they should explain the special circumstances that require a deviation from the Code of Ethics.
        \item The authors should make sure to preserve anonymity (e.g., if there is a special consideration due to laws or regulations in their jurisdiction).
    \end{itemize}


\item {\bf Broader impacts}
    \item[] Question: Does the paper discuss both potential positive societal impacts and negative societal impacts of the work performed?
    \item[] Answer: \answerNA{} % Replace by \answerYes{}, \answerNo{}, or \answerNA{}.
    \item[] Justification: This is a theoretical paper and does not discuss societal impacts.
    \item[] Guidelines:
    \begin{itemize}
        \item The answer NA means that there is no societal impact of the work performed.
        \item If the authors answer NA or No, they should explain why their work has no societal impact or why the paper does not address societal impact.
        \item Examples of negative societal impacts include potential malicious or unintended uses (e.g., disinformation, generating fake profiles, surveillance), fairness considerations (e.g., deployment of technologies that could make decisions that unfairly impact specific groups), privacy considerations, and security considerations.
        \item The conference expects that many papers will be foundational research and not tied to particular applications, let alone deployments. However, if there is a direct path to any negative applications, the authors should point it out. For example, it is legitimate to point out that an improvement in the quality of generative models could be used to generate deepfakes for disinformation. On the other hand, it is not needed to point out that a generic algorithm for optimizing neural networks could enable people to train models that generate Deepfakes faster.
        \item The authors should consider possible harms that could arise when the technology is being used as intended and functioning correctly, harms that could arise when the technology is being used as intended but gives incorrect results, and harms following from (intentional or unintentional) misuse of the technology.
        \item If there are negative societal impacts, the authors could also discuss possible mitigation strategies (e.g., gated release of models, providing defenses in addition to attacks, mechanisms for monitoring misuse, mechanisms to monitor how a system learns from feedback over time, improving the efficiency and accessibility of ML).
    \end{itemize}
    
\item {\bf Safeguards}
    \item[] Question: Does the paper describe safeguards that have been put in place for responsible release of data or models that have a high risk for misuse (e.g., pretrained language models, image generators, or scraped datasets)?
    \item[] Answer: \answerNA{} % Replace by \answerYes{}, \answerNo{}, or \answerNA{}.
    \item[] Justification: The paper does not release any data or models that have a high risk for misuse.
    \item[] Guidelines:
    \begin{itemize}
        \item The answer NA means that the paper poses no such risks.
        \item Released models that have a high risk for misuse or dual-use should be released with necessary safeguards to allow for controlled use of the model, for example by requiring that users adhere to usage guidelines or restrictions to access the model or implementing safety filters. 
        \item Datasets that have been scraped from the Internet could pose safety risks. The authors should describe how they avoided releasing unsafe images.
        \item We recognize that providing effective safeguards is challenging, and many papers do not require this, but we encourage authors to take this into account and make a best faith effort.
    \end{itemize}

\item {\bf Licenses for existing assets}
    \item[] Question: Are the creators or original owners of assets (e.g., code, data, models), used in the paper, properly credited and are the license and terms of use explicitly mentioned and properly respected?
    \item[] Answer: \answerNA{} % Replace by \answerYes{}, \answerNo{}, or \answerNA{}.
    \item[] Justification: The paper does not use existing assets.
    \item[] Guidelines:
    \begin{itemize}
        \item The answer NA means that the paper does not use existing assets.
        \item The authors should cite the original paper that produced the code package or dataset.
        \item The authors should state which version of the asset is used and, if possible, include a URL.
        \item The name of the license (e.g., CC-BY 4.0) should be included for each asset.
        \item For scraped data from a particular source (e.g., website), the copyright and terms of service of that source should be provided.
        \item If assets are released, the license, copyright information, and terms of use in the package should be provided. For popular datasets, \url{paperswithcode.com/datasets} has curated licenses for some datasets. Their licensing guide can help determine the license of a dataset.
        \item For existing datasets that are re-packaged, both the original license and the license of the derived asset (if it has changed) should be provided.
        \item If this information is not available online, the authors are encouraged to reach out to the asset's creators.
    \end{itemize}

\item {\bf New assets}
    \item[] Question: Are new assets introduced in the paper well documented and is the documentation provided alongside the assets?
    \item[] Answer: \answerNA{} % Replace by \answerYes{}, \answerNo{}, or \answerNA{}.
    \item[] Justification: The paper does not release new assets.
    \item[] Guidelines:
    \begin{itemize}
        \item The answer NA means that the paper does not release new assets.
        \item Researchers should communicate the details of the dataset/code/model as part of their submissions via structured templates. This includes details about training, license, limitations, etc. 
        \item The paper should discuss whether and how consent was obtained from people whose asset is used.
        \item At submission time, remember to anonymize your assets (if applicable). You can either create an anonymized URL or include an anonymized zip file.
    \end{itemize}

\item {\bf Crowdsourcing and research with human subjects}
    \item[] Question: For crowdsourcing experiments and research with human subjects, does the paper include the full text of instructions given to participants and screenshots, if applicable, as well as details about compensation (if any)? 
    \item[] Answer: \answerNA{} % Replace by \answerYes{}, \answerNo{}, or \answerNA{}.
    \item[] Justification: The paper does not involve crowdsourcing nor research with human subjects.
    \item[] Guidelines:
    \begin{itemize}
        \item The answer NA means that the paper does not involve crowdsourcing nor research with human subjects.
        \item Including this information in the supplemental material is fine, but if the main contribution of the paper involves human subjects, then as much detail as possible should be included in the main paper. 
        \item According to the NeurIPS Code of Ethics, workers involved in data collection, curation, or other labor should be paid at least the minimum wage in the country of the data collector. 
    \end{itemize}

\item {\bf Institutional review board (IRB) approvals or equivalent for research with human subjects}
    \item[] Question: Does the paper describe potential risks incurred by study participants, whether such risks were disclosed to the subjects, and whether Institutional Review Board (IRB) approvals (or an equivalent approval/review based on the requirements of your country or institution) were obtained?
    \item[] Answer: \answerNA{} % Replace by \answerYes{}, \answerNo{}, or \answerNA{}.
    \item[] Justification: The paper does not involve crowdsourcing nor research with human subjects.
    \item[] Guidelines:
    \begin{itemize}
        \item The answer NA means that the paper does not involve crowdsourcing nor research with human subjects.
        \item Depending on the country in which research is conducted, IRB approval (or equivalent) may be required for any human subjects research. If you obtained IRB approval, you should clearly state this in the paper. 
        \item We recognize that the procedures for this may vary significantly between institutions and locations, and we expect authors to adhere to the NeurIPS Code of Ethics and the guidelines for their institution. 
        \item For initial submissions, do not include any information that would break anonymity (if applicable), such as the institution conducting the review.
    \end{itemize}

\item {\bf Declaration of LLM usage}
    \item[] Question: Does the paper describe the usage of LLMs if it is an important, original, or non-standard component of the core methods in this research? Note that if the LLM is used only for writing, editing, or formatting purposes and does not impact the core methodology, scientific rigorousness, or originality of the research, declaration is not required.
    %this research? 
    \item[] Answer: \answerNA{} % Replace by \answerYes{}, \answerNo{}, or \answerNA{}.
    \item[] Justification: This paper does not involve LLMs as any important, original, or non-standard components.
    \item[] Guidelines:
    \begin{itemize}
        \item The answer NA means that the core method development in this research does not involve LLMs as any important, original, or non-standard components.
        \item Please refer to our LLM policy (\url{https://neurips.cc/Conferences/2025/LLM}) for what should or should not be described.
    \end{itemize}

\end{enumerate}

\end{document}
